\chapter{Funksies}
\setcounter{figure}{0}
\setcounter{subfigure}{0}

\section{Funksies in die real world}
Funksies is wiskundige boustene wat toepassings het in masjienontwerp, die voorspelling van natuurrampe, die mediese veld,
ekonomiese analise en vliegtuigontwerp. ’n Funksie het vir elke invoerwaarde net ’n enkele uitvoerwaarde. Dit is moontlik dat
’n funksie meer as een inset van verskillende veranderlikes kan hê, maar dan sal dit steeds net ’n enkele uitset hê. Ons gaan
egter nie in hierdie hoofstuk na sulke tipe funksies kyk nie.\par 
Een van die groot voordele van funksies is dat hulle ons toelaat om vergelykings te visualiseer deur middel van ’n grafiek. ’n
Grafiek is bloot ’n tekening van ’n funksie en dit word gebruik as ’n ander voorstellingswyse in plaas van ’n tabel met getalle.
In hierdie hoofstuk gaan ons leer hoe om funksies met reële getalle te skep en te verstaan; en hoe om grafieke te lees en te
teken.\par 
Funksies se toepassing omstrek groot wetenskap- en ingenieurs probleme tot alledaagse probleme. So, dit is nuttig om meer
te leer van funksies. ’n Funksie is altyd afhanklik van een of meer veranderlikes, soos tyd, afstand of ’n meer abstrakte entiteit.\par 

n Paar tipiese voorbeelde van funksies waarmee jy moontlik bekend is:\par 
\begin{itemize}[noitemsep]
\item Hoeveel geld jy het as ’n funksie van tyd. Hier is tyd die inset vir die funksie en die uitset is die bedrag geld. Jy sal op
enige oomblik net een bedrag geld hê. As jy verstaan hoe jou bedrag geld verander oor tyd, kan jy beplan hoe om jou
geld beter te spandeer. Besighede teken die grafiek van hulle geldsake oor tyd, sodat hulle kan sien wanneer hulle te
veel geld spandeer. Sulke waarnemings is nie altyd duidelik deur slegs na die getalle te kyk nie
\item Die temperatuur is ’n voorbeeld van ’n funksie met veelvuldige insette, insluitend die tyd van die dag, die seisoen, die
wolkbedekking, die wind, die plek en vele ander. Die belangrike ding om in te sien, is dat daar net een waarde vir
temperatuur is op ’n spesifieke plek, op ’n spesifieke tyd. As ons verstaan hoe die insette die temperatuur beïnvloed,
kan ons ons dag beter beplan.
\item Jou posisie is ’n funksie van tyd omdat jy nie op twee plekke op dieselfde tyd kan wees nie. Indien jy twee mense se
posisie as ’n funksie van tyd sou teken of stip (’plot’), sal die plek waar die lyne kruis, aandui waar die mense mekaar
ontmoet. Hierdie idee word gebruik in logistiek – ’n veld van Wiskunde wat probeer voorspel waar mense en items is,
hoofsaaklik vir besigheid.
\item Jou massa is ’n funksie van hoeveel jy eet en hoe baie oefening jy doen, maar elke persoon se liggaam hanteer die
insette anders en mens kan dan verskillende liggame voorstel as verskillende funksies.
\end{itemize}

\Definition{function}{A function is a mathematical relationship between two variables, where every input variable has one output.}
\Note{One output can have different input variables.}




\subsection*{Afhanklike en Onafhanklike Veranderlikes}
In functions, the $x$-value is known as the input or independent variable, because its value can be chosen freely. Die $y$-waarde
word gewoonlik bepaal deur een of ander verband gebaseer op ’n gegewe of gekose $x$-waarde. Ons noem die $x$-waarde die onafhanklike veranderlike, omdat die waarde vrylik gekies kan word. Die berekende $y$ -waarde is bekend as die afhanklike veranderlike, omdat die waarde afhanklik is van die gekose $x$-waarde.\par 

\subsection*{Definisieversameling en Waardeversameling}

Die definisieversameling (ook bekend as die gebied) van ’n verband is die stel ($x$-waardes) waarvoor daar te minste een ($y$-waarde) bestaan. 
\\Die waardeversameling (ook bekend as die terrein) is die stel $y$ waardes wat
bepaal kan word deur te minste een ($x$-waarde).\par 



\subsection*{Versamelkeurdernotasie}
Voorbeelde:
\\
\begin{table}[H]
\begin{tabular}{ |p{5cm} | p{8cm} | }
\hline
  $\{x: x \in \mathBB{R}, x > 0\}$ &  die stel van alle $x$ waardes, waar $x$  ’n reële getal groter as $0$ is.
\\ \hline
    $\{y: y \in \mathBB{N}, 3 < y \leq 5}$ & The set of all $y$-values such that $y$ is an element of the set of natural numbers, is greater than $3$ and less than or equal to $5$. 
\\ \hline
  $\{z: z \in \mathBB{Z}, z \leq 100}$ & The set of all $z$-values such that $z$ is an element of the set of integers and is less than or equal to $100$.  
\\ \hline
\end{tabular}
\end{table}
\subsection*{Intervalnotasie}
Hierdie notasie kan nie gebruik word om heelgetalle in ’n interval te beskryf nie.
Voorbeelde:
\\
\begin{table}[H]
\begin{tabular}{ |p{5cm} | p{8cm} | }
\hline
  $(4;12)$ &  ’n ronde hakie beteken die getal word uitgesluit uit die interval. This interval indicates all real numbers greater than, but not equal to $4$ and less than, but not equal to $12$.
\\ \hline
 $(- \infty; -1)$ & Round brackets are always used for positive and negative infinity. 
\\ \hline
 $[1; 13)$ & ’n Reghoekige hakie beteken die getal word ingesluit by die interval. This interval indicates all real numbers greater than, and equal to $1$ and less than, but not equal to $13$.
\\ \hline
\end{tabular}
\end{table}

\subsection*{Function notation}
This is a very useful way to express a function. Instead of writing $y=2x+1$ we write $f(x) = 2x+1$. We say ``$f$ of $x$ is equal to $2x+1$''. Any letter can be used, e.g. $g(x)$, $h(x)$, $p(x)$ etc. 
\begin{enumerate}[noitemsep, label=\textbf{\arabic*}. ] 
 \item Determine the output value: ``find the value of the function for $x=-3$'' can be written as`` find $f(-3)$''.
\\Replace $x$ with $-3$: $f(-3)=2(-3)+1=-5 \therefore f(-3)=-5$. \\
This means that when $x=-3$ the value of the function $(y)$ $~=5$.
\item Determine the input value: ``find the value of $x$ that will give a $y$-value of $27$'' can be written as ``find $x$ if $f(x)=2x+1 = 27$''. \\
We write the following equation and solve for $x$: $2x+1 = 27 \therefore x=13$.\\
This means that when $x=13$ the value of the function $(y)$ $~=27$.

\end{enumerate}
\Note{$x$-value = input value\\
$y$-value = output/function value}

\subsection*{Representations of functions}
Functions can be expressed in many different ways for different purposes. 
\begin{enumerate}[noitemsep, label=\textbf{\arabic*}. ] 
 \item Words: the relationship between two variables is such that one is always $5$ less than the other.
\item Mapping diagram: 
\begin{figure}[H]
\begin{center}
\scalebox{1} % Change this value to rescale the drawing.
{
\begin{pspicture}(0,-1.1360937)(6.7240624,1.1360937)
\psframe[linewidth=0.04,dimen=outer](3.7046876,0.32890624)(2.6246874,-0.75109375)
\psline[linewidth=0.04cm,arrowsize=0.05291667cm 2.0,arrowlength=1.4,arrowinset=0.4]{->}(4.1446877,0.14890625)(5.0246873,0.40890625)
\psline[linewidth=0.04cm,arrowsize=0.05291667cm 2.0,arrowlength=1.4,arrowinset=0.4]{->}(4.1446877,-0.57109374)(5.0046873,-0.85109377)
\psline[linewidth=0.04cm,arrowsize=0.05291667cm 2.0,arrowlength=1.4,arrowinset=0.4]{->}(4.1046877,-0.21109375)(4.9846873,-0.21109375)
\psline[linewidth=0.04cm,arrowsize=0.05291667cm 2.0,arrowlength=1.4,arrowinset=0.4]{<-}(2.073937,-0.58032197)(1.2008146,-0.86255836)
\psline[linewidth=0.04cm,arrowsize=0.05291667cm 2.0,arrowlength=1.4,arrowinset=0.4]{<-}(2.055675,0.13944638)(1.1888499,0.39754358)
\psline[linewidth=0.04cm,arrowsize=0.05291667cm 2.0,arrowlength=1.4,arrowinset=0.4]{<-}(2.104793,-0.21942326)(1.2250762,-0.24174327)
% \usefont{T1}{ptm}{m}{n}
\rput(0.44671875,0.93890625){Input:}
% \usefont{T1}{ptm}{m}{n}
\rput(6.0242186,0.93890625){Output:}
\usefont{T1}{ptm}{m}{n}
\rput(0.82609373,0.41890624){$-3$}
\usefont{T1}{ptm}{m}{n}
\rput(0.7753125,-0.22109374){$0$}
\usefont{T1}{ptm}{m}{n}
\rput(0.82953125,-0.92109376){$5$}
\usefont{T1}{ptm}{m}{n}
\rput(3.1746874,-0.19609375){\large $-5$}
\usefont{T1}{ptm}{m}{n}
\rput(5.368125,0.45890626){$-8$}
\usefont{T1}{ptm}{m}{n}
\rput(5.3646874,-0.24109375){$-5$}
\usefont{T1}{ptm}{m}{n}
\rput(5.3553123,-0.98109376){$0$}
\end{pspicture} 
}
\end{center}
\end{figure}
\item Table: 

 \begin{table}[H]
\begin{center}
  \begin{tabular}{|c|c|c|c|}
   \hline
Input variable & $-3$&$0$&$5$
\\ \hline
Output variable &$-8$&$-5$&$0$
\\ \hline
  \end{tabular}
\end{center}
 \end{table}



\item Set of ordered number pairs: $(-3;-8), (0;-5), (5;0)$
\item Algebraic formula: $y = x-5$, $x \in \{-3; 0; 5\}$
\item Graph:
\begin{figure}[H]
\begin{center}
\scalebox{1} % Change this value to rescale the drawing.
{
\begin{pspicture}(0,-3.0284376)(6.52125,3.0684376)
\rput(3.0,-0.0284375){\psaxes[linewidth=0.04,arrowsize=0.05291667cm 2.0,arrowlength=1.4,arrowinset=0.4,labels=none,ticks=none,ticksize=0.10583333cm]{<->}(0,0)(-3,-3)(3,3)}
\usefont{T1}{ptm}{m}{n}
\rput(6.3754687,0.1615625){$x$}
\usefont{T1}{ptm}{m}{n}
\rput(3.3557813,2.9415624){$y$}
\psline[linewidth=0.04cm,arrowsize=0.05291667cm 2.0,arrowlength=1.4,arrowinset=0.4]{<->}(2.0,-2.0884376)(5.18,0.8915625)
\psdots[dotsize=0.16](2.98,-1.1684375)
\psdots[dotsize=0.16](4.16,-0.0284375)
\usefont{T1}{ptm}{m}{n}
\rput(4.1448436,0.2415625){$5$}
\usefont{T1}{ptm}{m}{n}
\rput(3.3,-1.2384375){$-5$}
\usefont{T1}{ptm}{m}{n}
\rput(2.9814062,-0.0584375){$O$}
\end{pspicture} 
}
\end{center}
\end{figure}
\end{enumerate}

\begin{exercises}{}
{
Write the following in set notation:
\begin{enumerate}[noitemsep, label=\textbf{\arabic*}. ] 
 \item $(- \infty; 7]$
\item $[13;4)$
\item $(35; \infty)$
\item $[\frac{3}{4}; 21)$
\item $[-\frac{1}{2}; \frac{1}{2}$
\item $(-3; \infty)$
\end{enumerate}

Write the following in interval notation:
\begin{enumerate}[noitemsep, label=\textbf{\arabic*}. ] 
\setcounter{enumi}{6}
 \item $\{p: p \in \mathBB{R}, p \leq 6\}$
 \item $\{k: k \in \mathBB{R}, -5 < k < 5}$
 \item $\{x: x \in \mathBB{R}, x > \frac{1}{5}\}$
 \item $\{z: z \in \mathBB{R}, 21 \leq z < 41\}$
\end{enumerate}

} 
\end{exercises}

\section{reguitlynfunksies in die vorm $y=mx+c$}


\subsection*{Plotting the graph}       
Funksies met die algemene vorm $y=mx+c$ word reguitlynfunksies genoem. In die vergelyking, $y=mx+c$, $m$ en $c$ konstantes en het verskillende invloede op die grafiek van die funksie. 

\begin {wex}{Plotting a straight line}
{
\begin{figure}[H]
\begin{center}
\begin{pspicture}(-6,-2)(6,4.6)
\psset{yunit=0.3}
\psaxes[Dy=3]{<->}(0,0)(-6,-8)(6,14)
\psplot[plotstyle=curve,arrows=<->]{-5}{5}{x 2 mul 3 add}
\psplot[plotstyle=dots,arrows=<->,plotpoints=9]{-4}{4}{x 2 mul 3 add}
\rput(0, 0){$0$}
\rput(0.3, 13.6){$y$}
\rput(6.3, 0.3){$x$}
\rput(4.5, 9){$f(x)=2x+3$}
%\psplot[showpoints=true](-5,-7)(-5,-7)(-4,-5)(-3,-3)(-2,-1)(-1,1)(-1.5,0)(0,3)(1,5)(2,7)(3,9)(4,11)(5,13)
\rput(-1.6,1.5){$-\frac{3}{2}$}
\end{pspicture}
% \caption{Graph of $f(x) = 2x + 3$}
% \label{fig:mf:g:straightline}
\end{center}
\end{figure}  
}
{
\westep{Step 1}
Steps needed here
}
\end{wex}

  

\subsection*{Investigating the effects of $m$ and $c$}

\begin{Investigation}{The effects of $m$ and $c$ on a straight line graph}

Op dieselfde assestelsel, trek die volgende grafieke:
\begin{enumerate}[noitemsep, label=\textbf{\arabic*}. ] 

    \item $a(x)=x-2$
    \item $b(x)=x-1$
    \item $c(x)=x$
    \item $d(x)=x+1$
    \item $e(x)=x+2$
    \end{enumerate}
Gebruik jou resultate om die invloed van verskillende waardes van $c$ op die resulterende grafiek af te lei.\\
\\

Op dieselfde assestelsel, trek die volgende grafieke:

    \begin{enumerate}[noitemsep, label=\textbf{\arabic*}. ] 
\setcounter{enumi}{5}
    \item $f(x)=-2x$
    \item $g(x)=-x$
    \item $j(x)=x$
\item $k(x)=2x$
    \end{enumerate}
Gebruik jou resultate om die invloed van verskillende waardes van $m$ op die resulterende grafiek af te lei.
\end{Investigation}


\begin{table}[htb]
\begin{center}
% \caption{Table summarising general shapes and positions of graphs of functions of the form $y=mx+q=c$.}
\label{tab:mf:graphs:summarystr10}
\begin{tabular}{|c|c|c|c|}\hline
& $m<0$&$m=0$ & $m>0$\\ \hline
$c>0$&
\begin{pspicture}(-1.2,-1.2)(1.2,1.2)
\psset{yunit=0.25,xunit=0.25}
\psaxes[arrows=<->,dx=0,Dx=10,dy=0,Dy=10](0,0)(-4,-4)(4,4)
\psplot[plotstyle=curve,arrows=<->]{-2.5}{2.5}{x neg 1 add}
\end{pspicture}

&
\begin{pspicture}(-1.2,-1.2)(1.2,1.2)
\psset{yunit=0.25,xunit=0.25}
\psaxes[arrows=<->,dx=0,Dx=10,dy=0,Dy=10](0,0)(-4,-4)(4,4)
\psplot[plotstyle=curve,arrows=<->]{-2.5}{2.5}{1.5}
\rput(-0.5, 2.1){\footnotesize$q$}
\end{pspicture}
&
\begin{pspicture}(-1.2,-1.2)(1.2,1.2)
\psset{yunit=0.25,xunit=0.25}
\psaxes[arrows=<->,dx=0,Dx=10,dy=0,Dy=10](0,0)(-4,-4)(4,4)
\psplot[plotstyle=curve,arrows=<->]{-2.5}{2.5}{x 1 add}
\end{pspicture}
\\\hline
$c=0$&
\begin{pspicture}(-1.2,-1.2)(1.2,1.2)
\psset{yunit=0.25,xunit=0.25}
\psaxes[arrows=<->,dx=0,Dx=10,dy=0,Dy=10](0,0)(-4,-4)(4,4)
\psplot[plotstyle=curve,arrows=<->]{-2.5}{2.5}{x neg}
\end{pspicture}

&
---
&

\begin{pspicture}(-1.2,-1.2)(1.2,1.2)
\psset{yunit=0.25,xunit=0.25}
\psaxes[arrows=<->,dx=0,Dx=10,dy=0,Dy=10](0,0)(-4,-4)(4,4)
\psplot[plotstyle=curve,arrows=<->]{-2.5}{2.5}{x}
\end{pspicture}
\\ \hline
$c<0$
&

\begin{pspicture}(-1.2,-1.2)(1.2,1.2)
\psset{yunit=0.25,xunit=0.25}
\psaxes[arrows=<->,dx=0,Dx=10,dy=0,Dy=10](0,0)(-4,-4)(4,4)
\psplot[plotstyle=curve,arrows=<->]{-2.5}{2.5}{x neg 1 sub}
\end{pspicture}
&
\begin{pspicture}(-1.2,-1.2)(1.2,1.2)
\psset{yunit=0.25,xunit=0.25}
\psaxes[arrows=<->,dx=0,Dx=10,dy=0,Dy=10](0,0)(-4,-4)(4,4)
\psplot[plotstyle=curve,arrows=<->]{-2.5}{2.5}{1.5 neg}
\rput(-0.5, -2.1){\footnotesize$q$}
\end{pspicture}
&
\begin{pspicture}(-1.2,-1.2)(1.2,1.2)
\psset{yunit=0.25,xunit=0.25}
\psaxes[arrows=<->,dx=0,Dx=10,dy=0,Dy=10](0,0)(-4,-4)(4,4)
\psplot[plotstyle=curve,arrows=<->]{-2.5}{2.5}{x 1 sub}

\end{pspicture}
\\\hline
\end{tabular}
\end{center}
\end{table}

\textbf{The effect of $m$}\\
Jy behoort te vind dat die waarde van $m$ die helling van die grafiek beïnvloed. Soos $m$ vermeerder, vermeerder die
helling van die grafiek ook. Indien $m>0$ sal die grafiek vermeerder van links na regs (opwaartse helling). Indien $m<0$ sal die grafiek verminder van links na regs (afwaartse helling). Dit is hoekom daar na, $m$  verwys word as die helling of die gradiënt van ’n reguitlynfunksie.\par 

\textbf{The effect of $c$}\\
Jy behoort ook te vind dat die waarde van $c$ die punt bepaal waar die grafiek die $y$-as sny. Om hierdie rede, staan $c$ bekend as die y-afsnit. If $c$ increases the graph shifts vertically upwards. If $c$ decreases, the graph shifts vertically downwards.\par 

% These different properties are summarised in Table 1.7.\par 

\subsection*{Discovering the characteristics} 
Die algemene vorm van a reguitlynfunksie is:  $y=mx + c$. 
\subsection*{Definisieversameling en Waardeversameling}
\nopagebreak
Vir $f(x)=mx+c$, is die definisieversameling $\{x:x\in \mathbb{R}\}$ omdat daar geen waarde is van  $x\in \mathbb{R}$ waarvoor $f(x)$ ongedefinieërd is nie.\par 
Die waardeversameling van $f(x)=mx+c$ is ook $\{f(x):f(x)\in \mathbb{R}\}$ omdat daar geen waarde is van  $f(x)\in \mathbb{R}$ waarvoor $f(x)$ ongedefinieërd is nie.\par 
\par 

\subsection*{Afsnitte}
\textbf{Die $y$-afsnitte word as volg bereken:}\\
Every point on the $y$-axis has an $x$-coordinate of $0$. Therefore, to calculate the $y$-intercept, left $x=0$.\par
Byvoorbeeld, die $y$-afsnit van $g(x)=x-1$ word bepaal deur $x=0$ te stel en dan op te los:\par 

\begin{equation*}
\begin{array}{ccl}\hfill g(x)& =& x-1\hfill \\
\hfill g(0) &=& 0-1\hfill \\
& =& -1\hfill 
\end{array}
\end{equation*}
This gives the point $(0;-1)$.\par

\textbf{Die $x$-afsnitte word as volg bereken:}\\
Every point on the $x$-axis has a $y$-coordinate of $0$. Therefore to calculate the $x$-intercept, let $y=0$. \par
Byvoorbeeld, die $x$-afsnit van $g(x)=x-1$ word gegee deur $y=0$ in te stel en dan op te los:\par 

\begin{equation*}
\begin{array}{ccl}\hfill g(x)& =& x-1\hfill \\
\hfill 0& =& x-1\hfill \\
\hfill x& =& 1\hfill 
\end{array}
\end{equation*}
This gives the point $(1;0)$.


\subsection*{Skets van Grafieke van die vorm $f(x)=mx+c$}


Om die grafieke van die vorm, $f(x)=mx+c$, te skets, moet ons die volgende drie eienskappe vind:\par 
\begin{enumerate}[noitemsep, label=\textbf{\arabic*}. ] 
\item die teken van $m$
\item $y$-afsnit
 and axis of symmetry\item $x$-afsnit
\end{enumerate}
Slegs twee punte word benodig om ’n reguitlyn te trek. Die maklikste punte is die $x$-afsnit (waar die lyn die x-as
sny) en die $y$-afsnit.\par 

\subsection*{Dual intercept method}
\begin{wex}{Sketching a straigh line graph using the dual intercept method}
{skets die grafiek van $g(x)=x-1$ using the dual intercept method}
{
\westep{Examine the standard form of the equation}
$m>0$. This means that the graph increases as $x$ increases.

\westep{Calculate intercepts}
Die $y$-afsnit word bepaal deur $x=0$ te stel; Daarom $g(0)=-1$. en is bereken as $(0;-1)$. \\

Die $x$-afsnit word bepaal deur $y=0$ te stel; daarom $x=1$. en is bereken as $(1;0)$. 

\westep{Plot points and draw graph}

\begin{center}
\begin{pspicture}(-4,-4)(4,3)
%\psgrid
\psset{yunit=0.75,xunit=0.75}
\psaxes[arrows=<->](0,0)(-5,-5)(5,4)
\psplot[plotstyle=curve,arrows=<->]{-4}{4}{x 1 sub}
\psdots(0,-1)(1,0)
\uput[r](0,-1.3){$(0;-1)$}
\uput[ul](1.3,0.3){$(1;0)$}
\rput(5.2,0.3){$x$}
\rput(0.3, 4.2){$y$}
\end{pspicture}
% \caption{Graph of the function $g(x)=x-1$}
% \label{fig:mf:g:sketchexamplestr}
\end{center}

Note: this graph is continuous and extends in both directions.       
}
\end{wex}

\subsection*{Gradient and $y$-intercept method}
We can draw a straight line graph of the form $y=mx+c$ using the gradient $m$ and the $y$-intercept $c$. \par We calculate the $y$-intercept by letting $x=0$. This gives us one point for drawing the graph and to calculate the second we use the gradient ($m$).\par

The gradient of a line is the measure of steepness. Steepness is determined by the ratio of vertical change to horizontal change:
\begin{equation*}
m = \dfrac{\mbox{\footnotesize change in $y$}}{\mbox{\footnotesize change in $x$}} = \dfrac{\mbox{\footnotesize vertical change}}{\mbox{\footnotesize horizontal change}}
\end{equation*}
For example: $y=\frac{3}{2}x-1$, therfore $m > 0$ and $y$ increases as $x$ increases
\begin{equation*}
 m = \dfrac{\mbox{\footnotesize change in $y$}}{\mbox{\footnotesize change in $x$}} = \dfrac{3\uparrow}{2\rightarrow} = \dfrac{-3\downarrow}{-2\leftarrow}
\end{equation*}
\begin{center}
\scalebox{1} % Change this value to rescale the drawing.
{
\begin{pspicture}(0,-4.7584376)(9.34125,4.7584376)
\rput(4.0,0.3215625){\psaxes[linewidth=0.04,labels=none,ticksize=0.2cm, arrows=<->](0,0)(-4,-5)(5,4)}
\psline[linewidth=0.04cm,dotsize=0.07055555cm 2.0]{*-*}(1.92,-3.6984375)(5.96,2.3815625)
\psarc[linewidth=0.04](4.49,2.3515625){0.49}{0.0}{180.0}
\psarc[linewidth=0.04](5.51,2.3315625){0.51}{0.0}{180.0}
\rput{178.47418}(6.9182076,-7.536676){\psarc[linewidth=0.04](3.509283,-3.7222767){0.49}{0.0}{180.0}}
\rput{178.47418}(4.8816133,-7.4152513){\psarc[linewidth=0.04](2.4901774,-3.675124){0.51}{0.0}{180.0}}
\rput{88.38591}(3.753576,-4.18539){\psarc[linewidth=0.04](4.029283,-0.1622766){0.49}{0.0}{180.0}}
\rput{273.13806}(6.9809837,0.9839488){\psarc[linewidth=0.04](4.010177,-3.195124){0.51}{0.0}{180.0}}
\rput{273.13806}(5.9635777,1.9092364){\psarc[linewidth=0.04](3.9901774,-2.195124){0.51}{0.0}{180.0}}
\rput{273.13806}(4.9839826,2.874464){\psarc[linewidth=0.04](4.010177,-1.1951239){0.51}{0.0}{180.0}}
\rput{88.38591}(4.694869,-3.1535814){\psarc[linewidth=0.04](3.969283,0.8377234){0.49}{0.0}{180.0}}
\rput{88.38591}(5.6944723,-2.1817489){\psarc[linewidth=0.04](3.969283,1.8377234){0.49}{0.0}{180.0}}
\usefont{T1}{ptm}{m}{n}
\rput(9.195469,0.6515625){$x$}
\usefont{T1}{ptm}{m}{n}
\rput(4.355781,4.6315627){$y$}
\usefont{T1}{ptm}{m}{n}
\rput(5.1275,3.1515625){$+2\rightarrow$}
\usefont{T1}{ptm}{m}{n}
\rput(2.9310937,0.9715625){$+3\uparrow$}
\usefont{T1}{ptm}{m}{n}
\rput(5.161406,-2.1884375){$-3\downarrow$}
\usefont{T1}{ptm}{m}{n}
\rput(2.9578125,-4.6084375){$-2\leftarrow$}
\end{pspicture} 
}
\end{center}

\begin{wex}{Sketching a straigh line graph using the gradient-intercept method}
{Sketch the graph of $y=\frac{1}{2}x-3$ using the gradient-intercept method}
{
\westep{Use the intercept}
The value of $c$ tells us the $y$-intercept, where the graph cuts the $y$-axis. \\
$c=-3$ which gives us the point $(0;-3)$.

\westep{Use the gradient}

\begin{equation*}
 m = \dfrac{\mbox{\footnotesize change in $y$}}{\mbox{\footnotesize change in $x$}} = \dfrac{1\uparrow}{2\rightarrow} = \dfrac{-1\downarrow}{-2\leftarrow}
\end{equation*}
We start at  $(0;-3)$. Move $1$ unit up and $2$ units to the right. This gives us the second point $(2;-2)$. \\
Or move $1$ unit down and $2$ units to the lift. This gives us the second point $(-1,-4)$.

\westep{Plot points and draw graph}

\begin{center}
\begin{pspicture}(-5,-5)(5,5)
%\psgrid
\psset{yunit=0.75,xunit=0.75}
\psaxes[arrows=<->](0,0)(-6,-6)(7,4)
\psplot[plotstyle=curve,arrows=<->]{-4}{7}{x .5 mul 3 sub}
\psline[linewidth=.7pt,arrows=->](0,-2)(2,-2)
\psline[linewidth=.7pt,arrows=->](0,-3)(0,-2)
\psdots(0,-3)(2,-2)
\uput[r](0,-3.3){$(0;-3)$}
\uput[ul](4,-2.5){$(2;-2)$}
\rput(7.4,0.3){$x$}
\rput(0.3, 4.2){$y$}
\rput(-0.2, -2.5){\footnotesize$1$}
\rput(1, -1.7){\footnotesize$2$}
\end{pspicture}
% \caption{Graph of the function $g(x)=x-1$}
% \label{fig:mf:g:sketchexamplestr}
\end{center}

Note: this graph is continuous and extends in both directions.       
}
\end{wex}

\Tip{Always write thefunction in the form $y=mx+c$ and take note of $m$. After plotting the graph, make sure that if $m>0$, the graph increases and that if $m<0$, the graph decreases.}


\begin{exercises}{}
{
\nopagebreak

 List the $x$- and $y$-intercepts for the following straight-line graphs. Indicate whether the graph is increasing or decreasing:
    \begin{enumerate}[noitemsep, label=\textbf{\arabic*}. ] 
    \item $y=x+1$
    \item $y=x-1$
    \item $y=2x-1$
    \item $y+1=2x$
\item $3y-2x=6$
\item$y=-3$
\item $x=3y$
\item $\frac{x}{2} - \frac{y}{3} = 1$
    \end{enumerate}


For the functions in the diagram below, give the equation and domain and range:
    \begin{enumerate}[noitemsep, label=\textbf{\arabic*}. ] 
\setcounter{enumi}{8}
    \item $a(x)$
%etc - other functions on graph go here
    \end{enumerate} 
\setcounter{subfigure}{0}
\begin{figure}[H]
\begin{center}
\scalebox{1} % Change this value to rescale the drawing.
{
\begin{pspicture}(0,-4.1467185)(9.519062,4.1867185)
\rput(4.0,-0.14671862){\psaxes[linewidth=0.03,tickstyle=bottom,labels=none,ticks=none,ticksize=0.08cm, arrows=<->](0,0)(-4,-4)(4,4)}
\psline[linewidth=0.04cm](2.78,1.7732813)(7.76,-0.9067186)
\usefont{T1}{ptm}{m}{n}
\rput(4.2023435,3.9832811){$y$}
\usefont{T1}{ptm}{m}{n}
\rput(8.26,-0.036718626){$x$}
\usefont{T1}{ptm}{m}{n}
\rput(4.4696875,1.3432813){$(0;3)$}
\usefont{T1}{ptm}{m}{n}
\rput(6.9896874,0.043281376){$(4;0)$}
\usefont{T1}{ptm}{m}{n}
\rput(3.2445312,1.9167186){$a(x)$}
\psline[linewidth=0.04cm](3.2,-3.6332815)(8.42,2.3067186)
\usefont{T1}{ptm}{m}{n}
\rput(4.6996875,-2.8167186){$(0;-6)$}
\usefont{T1}{ptm}{m}{n}
\rput(8.894531,2.4967186){$b(x)$}
\psline[linewidth=0.04cm](0.98,1.1267186)(7.88,1.1067187)
\usefont{T1}{ptm}{m}{n}
\rput(8.3,1.1567186){$c(x)$}
\psline[linewidth=0.04cm](7.4,-2.0532813)(1.08,1.4467186)
\usefont{T1}{ptm}{m}{n}
\rput(7.9,-2.1032813){$d(x)$}
\usefont{T1}{ptm}{m}{n}
\rput(4.2,0.1){$0$}
\psline[linewidth=0.04](4.7,-0.87328136)(4.58,-0.45328137)(5.0,-0.45328137)
\psline[linewidth=0.04](4.94,0.30671862)(4.82,0.7267186)(5.24,0.7267186)
\end{pspicture} 
}
\end{center}

\end{figure}  
            
\begin{enumerate}[noitemsep, label=\textbf{\arabic*}. ] 
\setcounter{enumi}{12}
\item Sketch the following functions on the same set of axes, using the dual intercept method. Clearly indicating the intercepts with the axes as well as the co-ordinates of the point of interception of the graphs: $x+2y-5=0$ and $3x-y-1=0$
\item On the same system of axes, draw the graphs of $f(x)=3-3x$ and $g(x)=\frac{1}{3}x+1$ using the gradient-intercept method.
\end{enumerate}

\par \raisebox{-0.2em}{\includegraphics[height=1em]{../icons/www.pdf}} Find the answers with the shortcodes:
\par \begin{tabular}[h]{cccccc}
(1.) lxT  &  (2.) lxb  &  (3.) lxj  & \end{tabular}
}
\end{exercises}
   

\section{paraboliese funksies in die vorm $y=a{x}^{2}+q$}
\subsection*{Plotting the graph}         
Die algemene vorm en posisie van die grafiek van die funksie in die vorm $y=a{x}^{2}+q$  wat ons ’n parabool noem. In the equation $y=a{x}^{2}+q$, $a$ and $q$ are constant and have different effects on the parabola. 

\begin{wex}{Plotting a quadratic function}
{
\begin{equation*}
 y = f(x) = x^{2}
\end{equation*}

Voltooi die volgende tabel van funksiewaardes vir $f(x)=x^{2}$ en stel die waardes op dieselfde assestelsel:
\\
\begin{center}
\begin{tabular}{|c|c|c|c|c|c|c|c|}
\hline
  $x$ &  $-3$ & $-2$ & $-1$ & $0$ & $1$ & $2$ & $3$
\\ \hline
 $f(x)$& $9$ &&&&&&
\\ \hline
\end{tabular}
\end{center}
\vspace{10pt}
\begin{enumerate}[noitemsep, label=\textbf{\arabic*}. ] 
 \item Join the points with a smooth curve.
\item The domain of $f$ is $x \in \mathBB{R}$. Determine the range.
\item About which line is $f$ symmetrical?
\item From your graph, determine the value of $f(x) = \frac{25}{4}$. \\Confirm your answer graphically.
\item Write down the point where the graph cuts the axes.
\end{enumerate}
}
{
\westep{Substitute values into the equation}
\begin{equation*}
 \begin{array}{cclcc}
  f(x) &=& x^{2} & &\\
 f(-3) &=& (-3)^{2} &=& 9 \\ 
 f(-2) &=& (-2)^{2} &=& 4 \\
 f(-1) &=& (-1)^{2} &=& 1 \\
f(0) &=& 0^{2} &= &0 \\
f(1) &=& (1)^{2} &= &1 \\ 
f(2) &=& (2)^{2} &= &4 \\
f(3) &=& (3)^{2} &= &9
 \end{array}
\end{equation*}
\\
\\
\begin{center}
\begin{tabular}{|c|c|c|c|c|c|c|c|}
\hline
  $x$ &  $-3$ & $-2$ & $-1$ & $0$ & $1$ & $2$ & $3$
\\ \hline
 $f(x)$& $9$ &$4$&$1$&$0$&$1$&$4$&$9$
\\ \hline
\end{tabular}
\end{center}

\westep{Plot the points and join with a smooth curve}
From the table, we get the following points: \\
$(-3;9)$, $(-2;4)$, $(-1;1)$, $(0;0)$, $(1;1)$, $(2;4)$, $(3;9)$ \\
and the following graph: 
\begin{figure}[H]
\begin{center}
\begin{pspicture}(-5,-1)(5,5)
%\psgrid-\infty ;q
\psaxes[arrows=<->,dy=0.5](0,0)(-5,-1)(5,5)
\psset{yunit=0.5}
\psplot[plotstyle=curve,arrows=<->]{-3}{3}{x 2 exp}
\psdots(-2.5, 6.25)(2.5,6.25)
\rput(-2.2, 6.3){$B$}
\rput(2.2, 6.3){$A$}
\rput(0.1, 0.3){$0$}
\end{pspicture}
% \caption{Graph of $f(x)=x^2-1$.}
\label{fig:mf:g:parabola10}
\end{center}
\end{figure}    

\westep{Determine domain and range}
Domain: $x \in \mathBB{R}$\\
From the graph we see that for all values of $x$, $y$ is greater than or equal to $0$.\\
Range: $y \in [0; \infty)$

\westep{Find the axis of symmetry}
$f$ is symmetrical about the $y$-axis. Therefore the axis of symmetry of $f$ is the line $x=0$. 

\westep{Determine the $x$-value}
\begin{equation*}
 \begin{array}{ccl}
f(x) &=& \frac{25}{4} \\
\therfore \frac{25}{4} &=& x^{2} \\
x &=& \pm \frac{5}{2} 
\end{array}
\end{equation*}
See points $A$ and $B$ on the graph.

\westep{Determine the intercept}
The function $f(x)$ intercepts the axes at the origin $(0;0)$. \\
We notice that as the value of $x$ increases from $-\infty$ to $0$, $f(x)$ decreases. At the turning point $(0;0)$, $f(x) = 0$. As the value of $x$ increases from $0$ to $\infty$, $f(x)$ increases.
}
\end{wex}




  

\subsection*{Investigating the effects of $a$ and $q$ }
\begin{Investigation}{The effects of $a$ and $q$ on a parabolic graph}
Complete the table and plot the following graphs on the same system of axes:
    \begin{enumerate}[noitemsep, label=\textbf{\arabic*}. ] 
  \item $a(x)={x}^{2}-2$
    \item $b(x)={x}^{2}-1$
    \item $c(x)={x}^{2}$
    \item $d(x)={x}^{2}+1$
    \item $e(x)={x}^{2}+2$
        \end{enumerate}

\begin{table}[H]
% \begin{table}[H]
% \\ '' '0'
\begin{center}

\noindent

\begin{tabular}{|l|l|l|l|l|l|}\hline
 $x$&
$-2$&
$-1$&
$0$&
$1$&
$2$
\\ \hline


$a(x)$
&
&
&
&
&
\\ \hline

$b(x)$
&
&
&
&
&
\\ \hline

$c(x)$
&
&
&
&
&
\\ \hline

$d(x)$
&
&
&
&
&
\\ \hline

$e(x)$
&
&
&
&
&
\\ \hline

\end{tabular}
\end{center}
% \begin{center}{\small\bfseries Table 1.8}\end{center}
% \begin{caption}{\small\bfseries Table 1.8}\end{caption}
\end{table}
Use your results to deduce the effect of $q$.
\\
Complete the table and plot the following graphs on the same system of axes:
    \begin{enumerate}[noitemsep, label=\textbf{\arabic*}. ] 
\setcounter{enumi}{5}
  \item $f(x)=-2{x}^{2}$
    \item $g(x)=-{x}^{2}$
    \item $h(x)={x}^{2}$
    \item $i(x)=2{x}^{2}$
    \end{enumerate}
\begin{table}[H]
\begin{center}
\begin{tabular}{|l|l|l|l|l|l|}\hline
$x$&
$-2$&
$-1$&
$0$&
$1$&
$2$
\\ \hline

$f(x)$&
&
&
&
&
\\ \hline

$g(x)$&
&
&
&
&
\\ \hline


$h(x)$
&
&
&
&
&
\\ \hline

$i(x)$
&
&
&
&
&
\\ \hline

\end{tabular}
\end{center}
% \begin{center}{\small\bfseries Table 1.8}\end{center}
% \begin{caption}{\small\bfseries Table 1.8}\end{caption}
\end{table}

Use your results to deduce the effect of $a$.

\end{Investigation}

\begin{table}[H]
\begin{center}
% \caption{Table summarising general shapes and positions of graphs of functions of the form $y=mx+q=c$.}
\label{tab:mf:graphs:summarystr10}
\begin{tabular}{|c|c|c|}
\hline
 & $a<0$ & $a>0$
\\ \hline
$q>0$&
\begin{pspicture}(-1.2,-1.2)(1.2,1.2)
\psset{yunit=0.25,xunit=0.25}
\psaxes[arrows=<->,dx=0,Dx=10,dy=0,Dy=10](0,0)(-4,-4)(4,4)
\psplot[plotstyle=curve,arrows=<->]{-1.6}{1.6}{x 2 exp neg 1 add}
\end{pspicture}

&

\begin{pspicture}(-1.2,-1.2)(1.2,1.2)
\psset{yunit=0.25,xunit=0.25}
\psaxes[arrows=<->,dx=0,Dx=10,dy=0,Dy=10](0,0)(-4,-4)(4,4)
\psplot[plotstyle=curve,arrows=<->]{-1.6}{1.6}{x 2 exp 1 add}
\end{pspicture}
\\\hline
$q=0$&
\begin{pspicture}(-1.2,-1.2)(1.2,1.2)
\psset{yunit=0.25,xunit=0.25}
\psaxes[arrows=<->,dx=0,Dx=10,dy=0,Dy=10](0,0)(-4,-4)(4,4)
\psplot[plotstyle=curve,arrows=<->]{-1.6}{1.6}{x 2 exp neg}
\end{pspicture}
&
\begin{pspicture}(-1.2,-1.2)(1.2,1.2)
\psset{yunit=0.25,xunit=0.25}
\psaxes[arrows=<->,dx=0,Dx=10,dy=0,Dy=10](0,0)(-4,-4)(4,4)
\psplot[plotstyle=curve,arrows=<->]{-1.6}{1.6}{x 2 exp }
\end{pspicture}

\\ \hline
$q<0$
&

\begin{pspicture}(-1.2,-1.2)(1.2,1.2)
\psset{yunit=0.25,xunit=0.25}
\psaxes[arrows=<->,dx=0,Dx=10,dy=0,Dy=10](0,0)(-4,-4)(4,4)
\psplot[plotstyle=curve,arrows=<->]{-1.6}{1.6}{x 2 exp neg 1 sub}
\end{pspicture}
&

\begin{pspicture}(-1.2,-1.2)(1.2,1.2)
\psset{yunit=0.25,xunit=0.25}
\psaxes[arrows=<->,dx=0,Dx=10,dy=0,Dy=10](0,0)(-4,-4)(4,4)
\psplot[plotstyle=curve,arrows=<->]{-1.6}{1.6}{x 2 exp 1 sub}
\end{pspicture}
\\\hline
\end{tabular}
\end{center}
\end{table}

\textbf{Die effek $q$}
\\
Die waarde van $q$ is called a vertical shift because all points are moved the same distance in the same direction (because it slides the entire graph up or down). 
\begin{itemize}
\item Indien $q>0$, then the graph of $f(x)$ is shifted vertically upwards by $q$ units.Die draaipunt van $f(x)$ is bokant die $y$-as.
\item Indien $q<0$, then the graph of $f(x)$ is shifted vertically downwards by $q$ units. Die draaipunt van $f(x)$ is onderkant die $y$-as.
\end{itemize}
\textbf{The effect of $a$}
\\
Die waarde van $a$ bepaal die vorm van die grafiek. 
\begin{itemize}
 \item Indien $a>0$, sal die grafiek $f(x)$  ``glimlag'' and has a minimum turning point at $(0;q)$.\\
The graph of $f(x)$ is stretched vertically upwards; as $a$ gets larger, the graph gets narrower.
\\For $0<a<1$, $a$ is a fraction, which makes the graph of $f(x)$ get wider.
\item If $a<0$, sal die grafiek van $f(x)$ frons' and has a maximum turning point at $(0;q)$. 
\\The graph of $f(x)$ is stretched vertically downwards; as $a$ gets smaller, the graph gets narrower. \\
For $-1<a<0$, $a$ is a fraction, which makes the graph of $f(x)$ get wider.
\end{itemize}

\setcounter{subfigure}{0}
\begin{figure}[!ht]
\begin{center}
\begin{pspicture}(-4,-0.5)(6,2)
%\psgrid
\psset{yunit=0.5}
\psplot[plotstyle=curve,arrows=<->]{-2}{0}{x 1 add 2 exp}
\psplot[plotstyle=curve,arrows=<->]{3}{5}{x 4 sub 2 exp neg 1 add}
\psdots(-1.5,3)(-0.5,3)(3.5,3)(4.5,3)
\uput[d](-1,-0.5){$a>0$ ($a$ 'n positiewe glimlag)}
\uput[d](4,-0.5){$a<0$ ($a$ 'n negatiewe frons)}
\end{pspicture}
% \caption{Distinctive shape of graphs of a parabola if $a>0$ and $a<0$.}
\label{fig:mf:g:parabola10a}
\end{center}
\end{figure}   

This simulation allows you to visualise the effect of changing $a$ and $q$. Note that in this simulation $q = c$. Also an extra term $bx$ has been added in. You can leave $bx$ as $0$, or you can also see what effect this has on the graph.
\par 
\setcounter{subfigure}{0}
\begin{figure}[H] % horizontal\label{m39345*Ohms-Law}
\textnormal{Phet simulation for graphing}\vspace{.1in} \nopagebreak
\label{m39345*phet!!!underscore!!!sim}\label{m39345*phet-simulation}
\raisebox{-5 pt}{ \includegraphics[width=0.5cm]{col11306.imgs/summary_www.png}} { (Simulation:  MG10018 )}
\vspace{2pt}
\vspace{.1in}
\end{figure}       


\subsection*{Discovering the characterstics}
Die algemene vorm van funksies van die vorm  $y=ax^{2} + q$.
\subsection*{Definisieversameling en Waardeversameling}

Vir $f(x)=a{x}^{2}+q$, is die definisieversameling $\{x:x\in \mathbb{R}\}$ omdat daar nie ’n waarde is van $x\in \mathbb{R}$ waarvoor $f(x)$ ongedefinieërd is nie.\par 
\par 
Indien $a>0$ dan het ons:
\begin{equation*}
\begin{array}{cccl}\hfill {x}^{2}& \ge & 0\hfill & (\mbox{die kwadraat van ’n uitdrukking is altyd positief})\hfill \\
 \hfill a{x}^{2}& \ge & 0\hfill & (\mbox{since } a>0)\hfill \\
 \hfill a{x}^{2}+q& \ge & q\hfill & (\mbox{vermenigvuldiging met $q$ ’n positiewe getal, behou die volgorde van die ongelykheid}) \\
 \hfill \therefore f(x)& \ge & q\hfill & 
\end{array}
\end{equation*}
Dus indien $a>0$, is die waardeversameling gelyk aan $[q,\infty )$.\par 
Soortgelyk, kan ons aantoon indien  $a<0$ is die waardeversameling van $ (-\infty ,q]$. 

\begin{wex}{Domain and range of a parabola}
{If $g(x)={x}^{2}+2$, determine the domain and range of the function.}
{
\westep{Determine the domain}
die gebied is $\{x:x\in \mathbb{R}\}$ want daar is geen waarde van $x\in \mathbb{R}$ waarvoor $g(x)$ ongedefinieërd is nie.
\westep{Determine the range}

Die terrein van $g(x)$ kan as volg bereken word:

\begin{equation*}
\begin{array}{ccc}\hfill {x}^{2}& \ge & 0\hfill \\
 \hfill {x}^{2}+2& \ge & 2\hfill \\
 \hfill g(x)& \ge & 2\hfill 
\end{array}
\end{equation*}
Dus die waardeversameling is gelyk aan $\{g(x):g(x)\geq 2\}$.
}
\end{wex}



\subsection*{Afsnitte}
\textbf{Die $y$-afsnit}\\
Every point on the $y$-axis has an $x$-coordinate of $0$, therefore to calculate the $y$-intercept let $x=0$.\\

Byvoorbeeld, die $y$-afsnit van $g(x)={x}^{2}+2$ word verkry deur $x=0$ te stel, en dan:\par 

\begin{equation*}
\begin{array}{ccl}\hfill g(x)& =& {x}^{2}+2\hfill \\ 
\hfill {y}& =& {0}^{2}+2\hfill \\
 & =& 2\hfill 
\end{array}
\end{equation*}
This gives us the point $(0;2)$.\par

\textbf{Die $x$-afsnit}\\
Every point on the $x$-axis has an $y$-coordinate of $0$, therefore to calculate the $x$-intercept let $y=0$.\\

Byvoorbeeld, die $x$-afsnit van $g(x)={x}^{2}+2$ word verkry deur $y=0$te stel, en dan:\par

\begin{equation*}
\begin{array}{ccl}\hfill g(x)& =& {x}^{2}+2\hfill \\
 \hfill 0& =& x^{2}+2\hfill \\
 \hfill -2& =& x^{2}\hfill 
\end{array}
\end{equation*}
Hierdie antwoord is nie reëel nie. Daarom het die grafiek van $g(x)={x}^{2}+2$ geen $x$-afsnitte nie. 

\subsection*{Draaipunte}

Die draaipunte van funksies van die vorm $f(x)=ax^{2}+q$ word gegee deur na die waardeversameling van
die funksie te kyk. 
\begin{itemize}
 \item Indien $a>0$,die grafiek van $f(x)$ is 'n ``glimlag'' en het 'n (minimum) draaipunt by $(0;q)$.
\item Indien $a<0$,  die grafiek van $f(x)$ is 'n frons has n die (maksimum)
draaipunt by $(0;q)$.
\end{itemize}


\subsection*{Simmetrie-asse}

Daar is een simmetrie-as vir die funksie met die vorm $f(x)=ax^{2}+q$ en dit gaan deur die draaipunt. Omdat die draaipunt op die $y$-as lê, waar $x=0$. 

\subsection*{Sketching graphs of the form $f(x)=ax^{2}+q$}

Om ’n grafiek te skets van die vorm, $f(x)=a^{2}+q$ het ons vyf eienskappe nodig:
\begin{enumerate}[noitemsep, label=\textbf{\arabic*}. ] 
\item die teken van $a$
\item $y$-afsnit
\item $x$-afsnit
\item draaipunte

\end{enumerate}

\begin{wex}
 {Sketse van Parabole}
{trek 'n grafiek van $g(x)=-\frac{1}{2}x^{2}-3$.  Merk die afsnitte, draaipunt en die simmetrie-as.}
{
\westep{Examine the standard form of the equation}
Eerstens sien ons dat $a<0$.  Dit beteken dat die grafiek ’n maksimum draaipunt het.
\westep{Calculate the intercepts}
Die $y$-afsnit word bepaal deur $x=0$ te stel:
\begin{equation*}
\begin{array}{ccl}\hfill {y}& =& -\frac{1}{2}(0)^{2}-3\hfill \\
 h(0) =& \frac{1}{0} &= \mbox{ongedefinieërd}\vspace{5pt} \\ 
 & =& -\frac{1}{2}(0)-3\hfill \\
 & =& -3\hfill 
\end{array}
\end{equation*}
Die koördinate van die $y$ -afsnit is dan $(0; -3)$.\\

Die $x$-afsnit word bepaal deur $y=0$ te stel:
\begin{equation*}
\begin{array}{ccl}\hfill 0& =& -\frac{1}2x^{2}-3\hfill \\ 
\hfill 3& =& -\frac{1}2x^{2}\hfill \\
 \hfill -3(2)& =& x^{2}\hfill \\
\hfill -6& =& x^{2}\hfill \\
 h(0) &= \frac{1}{0} =& \mbox{ongedefinieërd}\vspace{5pt} \\ 
\end{array}
\end{equation*}
Die oplossing van die vergelyking is nie reëel nie. Daarom is daar geen $x$-afsnitte nie.
\westep{Bepaal die  die draaipunt}
Die draaipunt is gelyk aan die y-afsnit wat ons bereken het as $(0;-3)$.
\westep{Plot the points and sketch the graph}
% \begin{figure}[!ht]
\begin{center}
\begin{pspicture}(-5,-5)(5,1)
%\psgrid
\psset{yunit=0.75,xunit=0.75}
\psaxes[arrows=<->](0,0)(-5,-7)(5,1)
\psplot[plotstyle=curve,arrows=<->]{-2.5}{2.5}{x 2 exp 0.5 mul neg 3 sub}
\psdots(0,-3)
\uput[r](0,-2.7){$(0;-3)$}
\rput(0.3, 0.8){$y$}
\rput (5.2, 0.2){$x$}
\end{pspicture}
% \caption{Graph of the function $f(x)=-\frac{1}{2}x^2-3$}
% % \label{fig:mf:g:sketchexamplepar10}
\end{center}
% \end{figure}
\\
Die
definisievesameling is alle reële getalle en die waardeversameling is $(- \infty; -3]$. 
Ons weet dat die y -as die simmetrie-as is waar $x=0$.
}

\end{wex}


\begin{wex}
{Sketse van Parabole}
{Skets die grafiek van $y={3x}^{2}+5$. Merk die afsnitte, draaipunt en die simmetrie-as.}
{
\westep{Examine the standard form of the equation}
Die teken van $a>0$. s positief. Die parabool sal dus ’n minimum-
draaipunt hê.
\westep{Bereken die afsnit punte}
Die $y$-afsnit word bepaal deur $x=0$te stel:
\begin{equation*}
\begin{array}{ccl}\hfill y& =& 3x^{2}+5\hfill \\
 \hfill y& =& 3(0)^{2}+5\hfill \\
 \hfill & =& 5\hfill 
\end{array}
\end{equation*}
Die koördinate van die $y$-as is dan $(0;5)$.
Die $x$-afsnit word bepaal deur $y=0$ te stel:
\begin{equation*}
\begin{array}{ccc}\hfill y& =& 3x^{2}+5\hfill \\
 \hfill 0& =& 3x^{2}+5\hfill \\
 \hfill x^{2}& =& -\frac{3}{5}\hfill 
\end{array}
\end{equation*}
Die oplossing van die vergelyking is nie reëel nie. Daarom is daar geen $x$ -afsnitte nie.
 $x$-intercepts.

\westep{bereken die draaipunte}
Die draaipunt is by (0, q). Vir hierdie funksie is q = 5, dus die
draaipunt is by $(0;5)$.

\westep{Plot the points and sketch the graph}
\begin{center}
\scalebox{1}
{
\psset{xunit=1.0cm,yunit=0.5cm,algebraic=true,dotstyle=o,dotsize=3pt 0,linewidth=0.8pt,arrowsize=3pt 2,arrowinset=0.25}
\begin{pspicture*}(-2,-2)(4,8)
\psaxes[xAxis=true,yAxis=true,Dx=1,Dy=2,ticksize=-2pt 0,subticks=2]{->}(0,0)(-2,0)(2,8)[,140] [,-40]
\rput{0}(0,5){\psplot{-2}{2}{x^2/2/0.17}}
% \rput[bl](0.56,5.34){$y = 3x^{2} + 5$}
\usefont{T1}{ptm}{m}{n}
\rput(0.3,7.6){$y$}
\usefont{T1}{ptm}{m}{n}
\rput(2.3,0.3){$x$}
\end{pspicture*}
}
\end{center}\\
Die definisievesameling is alle reële getalle $\{x:x \in \mathBB{R}\}$\\
die waardeversameling $\{y: y \geq 5, y \in \mathBB{R}\}$\\
Ons weet dat die y -as die simmetrie-as is $x=0$.
}

\end{wex}


\noindent
Die volgende video wys een manier om grafieke te trek. Let op dat die term "vertex" in die video gebruik word vir die draaipunt.
\setcounter{subfigure}{0}
\begin{figure}[H] % horizontal\label{m39345*parabola-1}
\textnormal{Khan academy video on graphing parabolas - 1}\vspace{.1in} \nopagebreak
\label{m39345*yt-media1}\label{m39345*yt-video1}
\raisebox{-5 pt}{ \includegraphics[width=0.5cm]{col11306.imgs/summary_www.png}} { (Video:  MG10019 )}
\vspace{2pt}
\vspace{.1in}
\end{figure} 

   
\begin{exercises}{}
{
\begin{enumerate}[noitemsep, label=\textbf{\arabic*}. ] 
\item  Wys dat indien $a<0$ is die waardeversameling van  $f(x)=ax^{2}+q$ is $\{f(x):f(x) \leq q \}$.
\item Skets die grafiek van die funksie $y=-x^{2}+4$ en toon al die afsnitte met die asse.
\item Twee parabole is geteken: $g:y=ax^{2}+p$ en $h:y=bx^{2}+q$.
\setcounter{subfigure}{0}
\begin{center}
\scalebox{1} % Change this value to rescale the drawing.
{
\begin{pspicture}(0,-2.61)(9.072187,2.61)
\psline[linewidth=0.04cm,arrowsize=0.05291667cm 2.0,arrowlength=1.4,arrowinset=0.4]{->}(0.0,-0.57)(9.04,-0.57 )
\psline[linewidth=0.04cm,arrowsize=0.05291667cm 2.0,arrowlength=1.4,arrowinset=0.4]{->}(4.54,-2.59)(4.56,2.59)
\psbezier[linewidth=0.04](1.8801522,2.249799)(2.0476475,1.0067352)(2.6697726,-0.42756912)(2.9090517 ,-0.81005025)(3.1483305,-1.1925315)(3.6508162,-1.83)(4.584004,-1.83)
\psbezier[linewidth=0.04](7.24,2.249799)(7.0725045,1.0067352)(6.4503794,-0.42756912)(6.2111006,-0.81005025)(5.971822,-1.1925315)(5.469336,-1.83)(4.536148 ,-1.83)
\psbezier[linewidth=0.04](1.2402416,-2.409998)(1.4477341,-1.215622)(2.2184203,0.16250414)(2.5148382,0.53000444)(2.811256,0.8975047)(3.4337332,1.5100052)(4.5897627,1.5100052)
\psbezier[linewidth=0.04](7.88,-2.409998 )(7.672508,-1.215622)(6.901821,0.16250414)(6.6054034,0.53000444)(6.3089857,0.8975047)(5.686508,1.5100052)(4.530479,1.5100052)
\usefont{T1}{ptm}{m}{it}
\rput(4.8129687,2.48){$y$}
\usefont{T1}{ptm}{m}{it}
\rput(8.908281 ,-0.3){$x$}
\usefont{T1}{ptm}{m}{n}
\rput(4.2504687,1.73){$23$}
\usefont{T1}{ptm}{m}{n}
\rput(4.269219,-2.07){$ -9$}
\usefont{T1}{ptm}{m}{n}
\rput(1.730625,0.39){ $(-4;7)$} 
\usefont{T1}{ptm}{m}{n}
\rput(7.240625,0.39){$(4;7)$}
\usefont{T1}{ptm}{m}{n}
\rput(6.4275,-0.81){ $3$}
\usefont{T1}{ptm}{m}{it}
\rput(7.369219,1.87){$g$}
\usefont{T1}{ptm}{m}{it} 
\rput{-0.1}(0.0029207817,0.01384337){\rput(7.9021873,-1.69){$h$}}
\end{pspicture} 
}       
\end{center}

    \begin{enumerate}[noitemsep, label=\textbf{\alph*}. ] 
    \item Vind die waardes van $a$ en $p$.
    \item Vind die waardes van $b$ en $q$.
    \item Vind die waardes van $x$ for which $g{x}\ge h{x}$.
    \item For what values of $x$ is $g$ toenemend?
    \end{enumerate}
\end{enumerate}

\par \raisebox{-0.2em}{\includegraphics[height=1em]{../icons/www.pdf}} Kry die oplossing met die kortkodes:
\par \begin{tabular}[h]{cccccc}
(1.) lxD  &  (2.) lxW  &  (3.) lxZ  & \end{tabular}
}
\end{exercises}   

\section{hiperboliese funksies van die vorm $y=\frac{a}{x}+q$}


\subsection*{Plotting the graph}  
Funksies van die vorm $y=\frac{a}{x}+q$ staan bekend as hiperboliese funksies. 

\begin{wex}
{Plotting a hyperbolic function}
{
\begin{equation*}
 y = h(x) = \frac{1}{x}
\end{equation*}

Complete the following table for $h(x) = \frac{1}{x}$ and plot the points on a system of axes.

\begin{table}[H]
\begin{center}
\begin{tabular}{|c|c|c|c|c|c|c|c|c|c|c|c|}
\hline
  $x$ &  $-3$ & $-2$ & $-1$ & $-\frac{1}{2}$ & $-\frac{1}{4}$ &$0$&$\frac{1}{4}$&$\frac{1}{2}$&$1$&$2$&$3$
\\ \hline
 $h(x)$& $-\frac{1}{3}$ &&&&&&&&&&
\\ \hline
\end{tabular}
\end{center}
\end{table}


\begin{enumerate}[noitemsep, label=\textbf{\arabic*}. ] 
 \item Join the points with smooth curves.
\item What happens if $x=0$?
\item Explain why the graph of $h(x)=\frac{1}{x}$ consists of two separate curves.
\item What happens to $h(x)$ as the value of $x$ becomes very small or very large?
\item The domain of $h(x)$ is $x \in \mathBB{R} - \{0\}$. Determine the range.
\item About which two lines is $h$ symmetrical?
\end{enumerate}
}
{
\westep{Substitute values into the equation}{
\begin{equation*}
 \begin{array}{cllll}
  h(x) &=& \frac{1}{x} & &\\
 h(-3) &=& \frac{1}{-3} &=& -\frac{1}{3} \vspace{5pt}\\ 
 h(-2) &=& \frac{1}{-2} &=& -\frac{1}{2}\vspace{5pt} \\ 
 h(-1) &=& \frac{1}{-1} &=& -1 \vspace{5pt}\\ 
 h(-\frac{1}{2}) &=& \dfrac{1}{-\frac{1}{2}} &=& -2 \vspace{5pt}\\ 
 h(-\frac{1}{4}) &=& \dfrac{1}{-\frac{1}{4}} &=& -4 \vspace{5pt}\\ 
 h(0) &=& \frac{1}{0} &=& \mbox{undefined}\vspace{5pt} \\ 
 h(\frac{1}{4}) &=& \dfrac{1}{\frac{1}{4}} &=& 4 \vspace{5pt}\\ 
 h(\frac{1}{2}) &=& \dfrac{1}{\frac{1}{2}} &=& 2 \vspace{5pt}\\ 
 h(1) &=& \frac{1}{1} &=& 1 \vspace{5pt}\\ 
 h(2) &=& \frac{1}{2} &=& \frac{1}{2} \vspace{5pt}\\ 
 h(3) &=& \frac{1}{3} &=& \frac{1}{3} \vspace{5pt}\\ 
 \end{array}
\end{equation*}

\begin{table}[H]
\begin{center}
\begin{tabular}{|c|c|c|c|c|c|c|c|c|c|c|c|}
\hline
  $x$ &  $-3$ & $-2$ & $-1$ & $-\frac{1}{2}$ & $-\frac{1}{4}$ &$0$&$\frac{1}{4}$&$\frac{1}{2}$&$1$&$2$&$3$
\\ \hline
 $h(x)$& $-\frac{1}{3}$ &$-\frac{1}{2}$&$-1$&$-2$&$-4$&undefined&$4$&$2$&$1$&$\frac{1}{2}$&$\frac{1}{3}$
\\ \hline
\end{tabular}
\end{center}
\end{table}
}
\westep{Plot the points and join with two smooth curves}
From the table we get the following points: $(-3; -\frac{1}{3})$, $(-2; -\frac{1}{2})$, $(-1;-1)$, $(-\frac{1}{2}; -2)$, $(-\frac{1}{4}; -4)$, $(\frac{1}{4}; -4)$, $(\frac{1}{2}; 2)$, $(1; 1)$, $(2; \frac{1}{2})$ and $(3; \frac{1}{3})$. \vspace{8pt} \\


% \setcounter{subfigure}{0}
% % \begin{figure}[tbp]
% \begin{center}
% \begin{pspicture}(-5,-2)(5,5)
% %\psgrid
% \psset{yunit=0.75,xunit=0.75}
% \psaxes[arrows=<->](0,0)(-5,-2.667)(5,6)
% \psplot[plotstyle=curve,arrows=<->]{-5}{-0.25}{x -1 exp }
% \psplot[plotstyle=curve,arrows=<->]{0.25}{5}{x -1 exp }
% % \psplot[linestyle=dotted,plotstyle=curve]{-5}{2}{x 2 add}
% \end{pspicture}
% % \caption{General shape and position of the graph of a function of the form $f(x)=\frac{a}{x} + q$.}
% \label{fig:mf:g:hyperbola10}
% \end{center}
% % \end{figure}      


For $x=0$ the function $h$ is undefined. This is called a discontinuity at $x=0$. \vspace{8pt} \\
$y=h(x) = \frac{1}{x}$ therefore we can write that $x \times y = 1$. Since the product of two positive numbers \textbf{and} the product of two negative numbers can be equal to $1$, the graph lies in the first and third quadrants.

\westep{Determine the asymptotes}
As the value of $x$ gets bigger, the value of $h(x)$ gets closer to, but does not equal $0$. This is called a horizontal asymptote, the line $y=0$. The same happens in the third quadrant; as $x$ gets smaller, $h(x)$ also approaches the $x$-axis asymptotically.\vspace{8pt} \\

We also notice that there is a vertical asymptote, the line $x=0$; as $x$ gets closer to $0$, $h(x)$ approaches the $y$-axis asymptotically.
\westep{Determine the range}
Domain: $x \in \mathBB{R} - \{0\}$\\
From the graph, we see that $y$ is defined for all values except $0$.\\
Range: $x \in \mathBB{R} - \{0\}$ 
\westep{Determine the lines of symmetry}
The graph of $h(x)$ has two axes of symmetry: the lines $y=x$ and $y=-x$. About these two lines, one half of the hyperbola is a mirror image of the other half. 
}
\end{wex}




\subsection*{Ondersoek the effects of $a$ and $q$ }
\begin{Investigation}{Investigating the effects of $a$ and $q$ on a hyperbolic graph}
 Op dieselfde assestelsel, trek die volgende grafieke:
    \begin{enumerate}[itemsep=3pt, label=\textbf{\arabic*}. ] 
    \item $a(x)=\dfrac{1}{x}-2$
    \item $b(x)=\dfrac{1}{x}-1$
    \item $c(x)=\dfrac{1}{x}$
    \item $d(x)=\dfrac{1}{x}+1$
    \item $e(x)=\dfrac{1}{x}+2$
\end{enumerate}
Gebruik jou resultate om die invloed van  $q$ af te lei.\par

Op dieselfde assestelsel, trek die volgende grafieke::
    \begin{enumerate}[itemsep=3pt, label=\textbf{\arabic*}. ] 
\setcounter{enumi}{5}
    \item $f(x)=\dfrac{-2}{x}$
    \item $g(x)=\dfrac{-1}{x}$
    \item $h(x)=\dfrac{1}{x}$
    \item $i(x)=\dfrac{2}{x}$
    \end{enumerate}
Gebruik jou resultate om die invloed van $a$ af te lei.
\end{Investigation}

\begin{table}[H]
\begin{center}
% \caption{Table summarising general shapes and positions of functions of the form $y=\frac{a}{x} + q$. The axes of symmetry are shown as dashed lines.}
\label{tab:mf:graphs:summaryhyp10}
\begin{tabular}{|c|c|c|}\hline
& $a<0$&$a>0$\\\hline
$q>0$&
\begin{pspicture}(-1.2,-1.2)(1.2,1.2)
%\psgrid
\psset{xunit=0.2,yunit=0.2}
\psaxes[arrows=<->,dx=0,Dx=10,dy=0,Dy=10](0,0)(-5,-5)(5,5)
\psplot[plotstyle=curve,arrows=<->]{-5}{-0.25}{x -1 exp neg 2 add}
\psplot[plotstyle=curve,arrows=<->]{0.25}{5}{x -1 exp neg 2 add}
\psplot[linestyle=dotted,plotstyle=curve]{-4}{4}{x neg 2 add}
\end{pspicture}
&

\begin{pspicture}(-1.2,-1.2)(1.2,1.2)
%\psgrid
\psset{xunit=0.2,yunit=0.2}
\psaxes[arrows=<->,dx=0,Dx=10,dy=0,Dy=10](0,0)(-5,-5)(5,5)
\psplot[plotstyle=curve,arrows=<->]{-5}{-0.25}{x -1 exp 2 add}
\psplot[plotstyle=curve,arrows=<->]{0.25}{5}{x -1 exp 2 add}
\psplot[linestyle=dotted,plotstyle=curve]{-4}{4}{x 2 add}
\end{pspicture}
\\\hline
$q=0$ & &
\\ \hline
$q<0$&
\begin{pspicture}(-1.2,-1.2)(1.2,1.2)
%\psgrid
\psset{xunit=0.2,yunit=0.2}
\psaxes[arrows=<->,dx=0,Dx=10,dy=0,Dy=10](0,0)(-5,-5)(5,5)
\psplot[plotstyle=curve,arrows=<->]{-5}{-0.25}{x -1 exp neg 2 sub}
\psplot[plotstyle=curve,arrows=<->]{0.25}{5}{x -1 exp neg 2 sub}
\psplot[linestyle=dotted,plotstyle=curve]{-2}{4}{x neg 2 sub}
\end{pspicture}
&

\begin{pspicture}(-1.2,-1.2)(1.2,1.2)
%\psgrid
\psset{xunit=0.2,yunit=0.2}
\psaxes[arrows=<->,dx=0,Dx=10,dy=0,Dy=10](0,0)(-5,-5)(5,5)
\psplot[plotstyle=curve,arrows=<->]{-5}{-0.25}{x -1 exp 2 sub}
\psplot[plotstyle=curve,arrows=<->]{0.25}{5}{x -1 exp 2 sub}
\psplot[linestyle=dotted,plotstyle=curve]{-4}{4}{x 2 sub}
\end{pspicture}
\\\hline
\end{tabular}
\end{center}
\end{table}

\textbf{The effect of $q$}\newline

Jy behoort ook te vind dat die waarde van $q$ bepaal of die grafiek bo of onder die $x$-as is die word genoem a vertical shift because all points are moved the same distance in the same direction (because it slides the entire graph up or down). 
\begin{itemize}
\item Vir $q>0$, then the graph of $h(x)$ is shifted vertically upwards by $q$ units. 
\item For $q<0$, then the graph of $h(x)$ is shifted vertically downwards by $q$ units.
\end{itemize}
The horizontal asymptote is the line $y=q$ and the vertical asymptote is always the $y$-axis, the line $x=0$.\par
\vspace{8pt}
\textbf{The effect of $a$}\newline
Jy behoort te vind dat die waarde van $a$ bepaal die vorm van die grafiek  en as die grafiek in die eeste en derde kwardrante of in die tweede
en vierde kwadrante van die Cartesiese vlak lê. 
\begin{itemize}
 \item Indien $a>0$, sal die grafiek $h(x)$ in die eeste en derde kwardrante lê. \\
For $a>1$, the graph of $h(x)$ will be further away from the axes than $y=\frac{1}{x}$.
\\For $0<a<1$, $a$ is a fraction and the graph will be closer to the axes than $y=\frac{1}{x}$. 
\item Indien $a<0$, sal die grafiek $h(x)$ in die tweede
en vierde kwadrante van die Cartesiese vlak lê.\\
For $a<-1$, the graph of $h(x)$ will be further away from the axes than $y=-\frac{1}{x}$.
\\For $-1<a<1$, $a$ is a fraction and the graph will be closer to the axes than $y=-\frac{1}{x}$. 
\end{itemize}



\subsection*{Discovering the characteristics}  
Die algemene vorms en posisies van funksies van die vorm $y=\frac{a}{x}+q$.

\subsection*{Definisieversameling en Waardeversameling}

Die funksie $y=\frac{a}{x}+q$, is ongedefiniëerd vir $x=0$. \\
Die definisieversameling is dus $\{x:x\in \mathbb{R},x\ne 0\}$.\par 
Ons kan sien dat $y=\frac{a}{x}+q$ herskryf kan word as:
\begin{equation*}
\begin{array}{ccl}\hfill y& =& \dfrac{a}{x}+q\hfill \vspace{4pt} \\
 \hfill y-q& =& \dfrac{a}{x}\hfill \\
 \hfill \mbox{Indien }x \ne  0 \mbox{dan}:(y-q)x& =& a\hfill \\
 \hfill x& =& \dfrac{a}{y-q}\hfill 
\end{array}
\end{equation*}
Dit wys dat die funksie ongedefiniëerd is by $y=q$. \\
Die waardeversameling van $f(x)=\frac{a}{x}+q$ is $\{f(x):f(x)\in (-\infty ;q)\cup (q;\infty )\}$.\par 

\begin{wex}{Definisieversameling en Waardeversameling of a hyperbola}
{If $g(x)=\frac{2}{x}+2$, determine the domain and range of the function.}
{
\westep{Bereken die definisieversameling}
The domain is $\{x:x\in \mathbb{R},x\ne 0\}$ omdat $g(x)$ ongedefiniëerd is by $x=0$.
\westep{Bereken die waardeversameling}
Ons sien dat  $g(x)$ ongedefiniëerd is by $y=2$.  Die waardeversamling is dus $\{g(x):g(x)\in (-\infty ;2)\cup (2;\infty )\}$.
}
\end{wex}


\subsection*{Afsnitte}

\textbf{The $y$-intercept} \\
Every point on the $y$-axis has an $x$-coordinate of $0$, therefore to calculate the $y$-intercept, let $x=0$.\\
Byvoorbeeld, vir die $y$-as van $g(x)=\frac{2}{x}+2$ word gegee deur $x=0$ te stel:
\begin{equation*}
\begin{array}{ccc}\hfill y& =& \dfrac{2}{x}+2\hfill \vspace{4pt}\\
 \hfill y& =& \dfrac{2}{0}+2\hfill 
\end{array}
\end{equation*}
Dit is ongedefiniëerd omdat ons deur nul deel. Daar is dus geen $y$-afsnit nie.\\
\\

\textbf{The $x$-intercept} \\
Every point on the $x$-axis has an $y$-coordinate of $0$, therefore to calculate the $x$-intercept, let $y=0$.\\
Byvoorbeeld, die  $x$ afsnit word bereken deur $g(x)=\frac{2}{x}+2$ word gegee deur $y=0$ te stel:
\begin{equation*}
\begin{array}{ccl}
\hfill y& =& \dfrac{a}{x}+q\hfill \vspace{4pt}\\
 \hfill 0& =& \dfrac{a}{x}+q\hfill \vspace{4pt} \\
 \hfill \dfrac{a}{x}& =& -q\hfill \\
 \hfill a& =& -q(x)\hfill \vspace{4pt}\\
 \hfill x& =& \dfrac{a}{-q}\hfill 
\end{array}
\end{equation*}
This gives us the point $(-1; 0)$.


\subsection*{Asimptote}

Daar is twee asimptote vir die funksies van die vorm $y=\frac{a}{x}+q$. \par 
Ons het gesien dat die funksie ongedefenieer was by $x=0$ en $y=q$ Dus is die asimtote $y=q$(die $y$-axis), en $x=0$. 

\subsection*{Axes of symmetry}
There are two lines about which a hyperbola is symmetrical: $y=ax+q$ and $y = -x +q$.


\subsection*{Skets die Grafieke van die vorm  $f(x)=\frac{a}{x}+q$}

Om grafieke van funksies van die vorm, $f(x)=\frac{a}{x}+q$, te skets, het ons vier eienskappe nodig.
\begin{enumerate}[noitemsep, label=\textbf{\arabic*}. ] 
\item teken van $a$
\item $y$-afsnitte
\item $x$-afsnitte
\item Asimptote
\end{enumerate}

\begin{wex}{Skets ’n Hiperbool}
{Skets die Grafiek $g(x)=\frac{2}{x}+2$. Merk die afsnitte en asimptote.}
{
\westep{Examine the standard form of the equation}
Ons sien dat $a>0$ daarom lê. die grafiek $g(x)$ in die eeste en derde kwardrante. 
\westep{Bereken die afsnitte}
For the $y$-intercept, let $x=0$:
\begin{equation*}
\begin{array}{ccl}
  g(x) = & \dfrac{2}{x} + 2  \vspace{4pt} \\
  g(0) = & \dfrac{2}{0} +2  
 \end{array}
\end{equation*}
Dit is ongedefiniëerd, Daar is geen $y$-afsnit nie. 
\\
For the $x$-intercept, let $y=0$:
\begin{equation*}
 \begin{array}{ccl}
  g(x) = &  \dfrac{2}{x} + 2 \vspace{4pt}\\
 0 = & \dfrac{2}{x} +2 \\
\dfrac{2}{x} = & -2 \\
\therefore x = &-1
 \end{array}
\end{equation*}

This gives the point $(-1;0)$


\westep{Determine the asymptotes}
Die asimptote kan dus gevind word by$y=2$ en $x=0$.

\westep{Skets die Grafieke}
\setcounter{subfigure}{0}
\begin{figure}[H]
\begin{center}
\begin{pspicture}(-5,-3)(5,6)
%\psgrid
\psset{yunit=0.75,xunit=0.75}
\psaxes[arrows=<->](0,0)(-5,-4)(5,7)
\psplot[plotstyle=curve,arrows=<->]{-5}{-0.4}{x -1 exp 2 mul 2 add}
\psplot[plotstyle=curve,arrows=<->]{0.4}{5}{x -1 exp 2 mul 2 add}
\psline[linestyle=dashed](-5,2)(5,2)
\end{pspicture}
% \caption{Graph of $g(x)=\frac{2}{x} + 2$.}
% \label{fig:mf:g:hyperbolasketchexample}
\end{center}
\end{figure} 
Domain: $x \in \mathBB{R} - \{0\}$\\
Range: $y \in \mathBB{R} - \{2\}$
}
\end{wex}




\begin{wex}
{Skets ’n Hiperbool}
{
Skets die grafiek van $y=\frac{-4}{x}+7$.}
{

\westep{Examine the standard form of the equation}
Ons sien dat $a<0$ therefore the graph lies in the second and fourth quadrants.
\westep{Calculate the intercepts}
Die $y$-afsnit is waar $x=0$:
\begin{equation*}
 \begin{array}{ccc}
 \hfill  y &= & \dfrac{-4}{x}+7 \vspace{4pt}\hfill \\
 \hfill &= & \dfrac {-4}{0} +7  \hfill \\

 \end{array}
\end{equation*}
Die funksie is ongedefiniëerd by $y$. Daar is geen $y$-afsnit vir grafieke van hierdie vorm nie. \\
Die $x$-afsnit is waar $y=0$:
\begin{equation*}
 \begin{array}{ccl}
 y &=&  \dfrac{-4}{x}+7\vspace{4pt}\\
 0 &=&  \dfrac{-4}{x}+7\vspace{4pt}\\ 
 \dfrac{-4}{x} &=& -7\vspace{4pt} \\
\therefore x &= &\dfrac{4}{7}
 \end{array}
\end{equation*}

Daar is dus een x-afsnit by $\Big(\dfrac{4}{7};0\Big)$


\westep{Determine the asymptotes}
Ons kyk na die gebied en die terrein om te bepaal waar die asimp-
tote lê. Van die gebied kan ons sien dat die funksie ongedefiniëerd
is wanneer $x=0$, Dus daar is een asimptoot by $x=0$.Die funksie is
ongedefiniëerd by y = q. Dus die tweede asimptoot is by $y=7$. 

\westep{Skets die grafiek}
\setcounter{subfigure}{0}
\begin{figure}[H]
\psset{xunit=0.5cm,yunit=0.5cm,algebraic=true,dotstyle=o,dotsize=3pt 0,linewidth=0.8pt,arrowsize=3pt 2,arrowinset=0.25}
\begin{pspicture*}(-10.85,-7.73)(10.36,13.74)
\psaxes[Axis=true,yAxis=true,Dx=2,Dy=2,ticksize=-2pt 0,subticks=2]{->}(0,0)(-10.85,-7.73)(10.36,13.74)[x,140] [y,-40]
\psplot[plotpoints=200]{-10.849663672784777}{10.361464563594}{-4/x+7}
\psplot[linestyle=dashed,dash=5pt 5pt]{-10.85}{10.36}{(--7-0*x)/1}
\rput[bl](-8.13,7.85){$y=-\frac{4}{x} + 7$}
\end{pspicture*}   
\end{figure}

Domain: $x \in \mathBB{R} - \{0\}$\\
Range: $y \in \mathBB{R} - \{7\}$\\
Axis of symmetry: $y=x+7$ and $y=-x+7$


}
\end{wex}

\begin{exercises}{}
{
Gebruik grafiekpapier en teken die grafiek van  $xy=-6$.
    \begin{enumerate}[noitemsep, label=\textbf{\arabic*}. ] 
    \item Lê die punt $(-2; 3)$  op die grafiek? Gee ’n rede vir jou antwoord.

    \item As die $x$-waarde van ‘n punt op die grafiek $0,25$ wat is die ooreenstemmende $y$-waarde?
    \item Wat gebeur met die $y$-waardes as die $x$-waardes baie groot word?
\item Calculate the coordinates of the points $P$origin to the graph.
\item Give the equations of the asymptotes.
    \item Met die lyn $y=-x$ as ’n lyn van simmetrie, watter punt is simmetries ten opsigte van $(-2; 3)$ ?
    \end{enumerate}
Skets die grafiek van  $h(x)=\frac{8}{x}$.
    \begin{enumerate}[noitemsep, label=\textbf{\arabic*}. ] 
\setcounter{enumi}{7}
    \item Hoe sal die grafiek $g(x)=\frac{8}{x}+3$ vergelyk met die grafiek van $h(x)=\frac{8}{x}$? Verduidelik jou antwoord.
    \item Skets die grafiek van $y=\frac{8}{x}+3$ op dieselfde assestelsel., showing asymptotes, axes of symmetry and the coordinates of one point on the graph.
    \end{enumerate}


\par \raisebox{-0.2em}{\includegraphics[height=1em]{../icons/www.pdf}} Find the answers with the shortcodes:
\par \begin{tabular}[h]{cccccc}
(1.) lxB  &  (2.) lxK  & \end{tabular}
}
\end{exercises}

\section{eksponensiële funksies van die vorm $y=ab^{x}+q$}

\subsection*{Plotting the graph}         
Funksies van die vorm $y=ab^{x}+q$ is bekend as eksponensiële funksies. In the equation $a$ aned $q$ are constants and have different effects on the function.

\begin{wex}{Plotting an exponential function}
 {
\begin{equation*} y=f(x) =b^{x} \mbox{ for } b>0 \mbox{ and } b \neq 1 \end{equation*}

Complete the following table for each of the functions and draw the graphs on the same system of axes:
$f(x)=2^{x}$, $g(x)=3^{x}$, $h(x)=5^{x}$.


\begin{table}[H]
\begin{center}
\begin{tabular}{|c|c|c|c|c|c|}
\hline
   &  $-2$ & $-1$ & $0$ & $1$ & $2$ 
\\ \hline
 $y=2^{x}$&  &&&&
\\ \hline
 $y=3^{x}$&  &&&&
\\ \hline
 $y=5^{x}$&  &&&&
\\ \hline

\end{tabular}
\end{center}
\end{table}

\begin{enumerate}[noitemsep, label=\textbf{\arabic*}. ] 
 \item At what point do these graphs intersect?
\item Explain why they do not cut the $x$-axis.
\item Give the domain and range of $h(x)$.
 \item As $x$ increases, does $h(x)$ increase or decrease?
\item Which of these graphs increases at the slowest rate?
\item For $y=k^{x}$ and $k>1$, the greater the value of $k$ the steeper the curve of the graph. True or false?
\end{enumerate}

Complete the following table for each of the functions and draw the graphs on the same system of axes:
$F(x) =(\frac{1}{2})^{x}$, $G(x) =(3)^{-x}$, $H(x) =(\frac{1}{5})^{x}$, 
\begin{table}[H]
\begin{center}
\begin{tabular}{|c|c|c|c|c|c|}
\hline
   &  $-2$ & $-1$ & $0$ & $1$ & $2$ 
\\ \hline
 $y=(\frac{1}{2})^{x}$&  &&&&
\\ \hline
$y=(\frac{1}{3})^{x}$&  &&&&
\\ \hline
$y=(\frac{1}{5})^{x}$&  &&&&
\\ \hline

\end{tabular}
\end{center}
\end{table}

\begin{enumerate}[noitemsep, label=\textbf{\arabic*}. ] 
 \item Give the $y$-intercepts for each function.
\item Describe the relationship between the graphs $f(x)$ and $F(x)$.
\item Describe the relationship between the graphs $g(x)$ and $G(x)$.
\item Give the domain and range of $H(x)$.
\item For $y=(\frac{1}{k})^{x}$ and $k>1$, the greater the value of $k$ the steeper the curve of the graph. True or false?
\item Give the equations of the asymptotes for each of the functions.
\end{enumerate}

}
{
\westep{Substitute values into the equations}
\begin{table}[H]
\begin{center}
\begin{tabular}{|c|c|c|c|c|c|}
\hline
   &  $-2$ & $-1$ & $0$ & $1$ & $2$ 
\\ \hline
 $f(x)=2^{x}$& $\frac{1}{4}$ &$\frac{1}{2}$&$1$&$2$&$4$
\\ \hline
 $g(x)=3^{x}$& $\frac{1}{9}$ &$\frac{1}{3}$&$1$&$3$&$9$
\\ \hline
 $g(x)=5^{x}$& $\frac{1}{25}$ &$\frac{1}{5}$&$1$&$5$&$25$
\\ \hline

\end{tabular}
\end{center}
\end{table}

\begin{table}[H]
\begin{center}
\begin{tabular}{|c|c|c|c|c|c|}
\hline
   &  $-2$ & $-1$ & $0$ & $1$ & $2$ 
\\ \hline
 $F(x)=(\frac{1}{2})^{x}$& $4$ &$2$&$1$&$\frac{1}{2}$&$\frac{1}{4}$
\\ \hline
$G(x)=(\frac{1}{3})^{x}$&  $9$&$3$&$1$&$\frac{1}{3}$&$\frac{1}{9}$
\\ \hline
$H(x)=(\frac{1}{5})^{x}$& $25$& $5$&$1$&$\frac{1}{5}$&$\frac{1}{25}$
\\ \hline

\end{tabular}
\end{center}
\end{table}
\westep{Plot the points and join with a smooth curve}
% Diagram needed here
\begin{enumerate}[noitemsep, label=\textbf{\arabic*}. ] 
\item We notice that all graphs pass through the point $(0;1)$. Any number with exponent $0$ is equal to $1$.
\item The graphs do not cut the $x$-axis because $0^{0}$ is undefined.
\item Domain: $x \in \mathBB{R}$\\
Range: $(0; \infty)$
\item As $x$ increases, $h(x)$ increases.
\item $f(x)=2^{x}$ increases at the slowest rate because it has the smallest base.
\item True: the greater the value of $k ~(k>1)$, the steeper the graph of $y=k^{x}$.
\end{enumerate}
% Another diagram needed here
\begin{enumerate}[noitemsep, label=\textbf{\arabic*}. ] 
\item The $y$-intercept is the point $(0; 1)$ for all the graphs. For any real number $z$, $z^{0}=1$.
\item $F(x)$ is the reflection of $f(x)$ about the $y$-axis. 
\item $G(x)$ is the reflection of $g(x)$ about the $y$-axis. 
\item  Domain: $x \in \mathBB{R}$\\
Range: $(0; \infty)$
\item True: the smaller the value of $k ~(k>1)$, the steeper the graph of $y=(\frac{1}{k})^{x}$.
\item The equation of the horizontal asymptote is $y=0$, the $x$-axis.
\end{enumerate}

}
\end{wex}

% \par 
% \setcounter{subfigure}{0}
% \begin{figure}[H]
% \begin{center}
% \begin{pspicture}(-5,-1)(5,4)
% %\psgrid
% \psset{yunit=0.75,xunit=0.75}
% \psaxes[arrows=<->](0,0)(-5,-1)(5,5)
% \psplot[plotstyle=curve,arrows=<->]{-5}{1.2}{2 x exp 2 add}
% \end{pspicture}
% % \caption{General shape and position of the graph of a function of the form $f(x)=ab^{x} + q$.}
% % \label{fig:mf:g:exponential10}
% \end{center}
% \end{figure}     

\subsection*{Investigating the effects of $a$ and $q$}
\begin{Investigation}{Investigating the effects of $a$ and $q$ on an exponential graph}
Op dieselfde assestelsel, skets die volgende grafieke ($b=2$, $a=1$ en $q$ verander):
\begin{enumerate}[noitemsep, label=\textbf{\alph*}. ] 
\item $a(x)=2^{x}-2$
\item $b(x)=2^{x}-1$
\item $c(x)=2^{x}$
\item $d(x)=2^{x}+1$
\item $e(x)=2^{x}+2$
\end{enumerate}

\begin{table}[H]
\begin{center}
\begin{tabular}{|l|c|c|c|c|c|}
\hline
   &  $-2$ & $-1$ & $0$ & $1$ & $2$ 
\\ \hline
$a(x)=2^{x}-2$&  &&&&
\\ \hline
 $b(x)=2^{x}-1$&  &&&&
\\ \hline
$c(x)=2^{x}$&  &&&&
\\ \hline
$d(x)=2^{x}+1$&  &&&&
\\ \hline
$e(x)=2^{x}+2$&  &&&&
\\ \hline
\end{tabular}
\end{center}
\end{table}
Gebruik jou antwoorde om ’n gevolgtrekking ten opsigte van die invloed van $q$ te maak.
\\

Op dieselfde assestelsel, skets die volgende grafieke ($b=2$, $q=0$ en $a$ verander):
\begin{enumerate}[noitemsep, label=\textbf{\arabic*}. ] 
\item $f(x)=2^{x}$
\item $g(x)=2.2^{x}$
\item $h(x)=-2^{x}$
\item $i(x)=-2.2^{x}$
\end{enumerate}

\begin{table}[H]
\begin{center}
\begin{tabular}{|l|c|c|c|c|c|}
\hline
   &  $-2$ & $-1$ & $0$ & $1$ & $2$ 
\\ \hline
$f(x)=2^{x}$&  &&&&
\\ \hline
$g(x)=2.2^{x}$&  &&&&
\\ \hline
 $h(x)=-2^{x}$&  &&&&
\\ \hline
$i(x)=-2.2^{x}$&  &&&&
\\ \hline

\end{tabular}
\end{center}
\end{table}
Gebruik jou antwoorde om ’n gevolgtrekking ten opsigte van die invloed van $a$ te maak.
\end{Investigation}


\begin{table}[H]
\begin{center}
% \caption{Table summarising general shapes and positions of functions of the form $y=ab^{x} + q$.}
% \label{tab:mf:graphs:summaryexp10}
\begin{tabular}{|c|c|c|}\hline
& $a<0$&$a>0$\\\hline
$q>0$&


\begin{pspicture}(-1.2,-1.2)(1.2,1.2)
%\psgrid
\psset{xunit=0.2,yunit=0.2}
\psaxes[arrows=<->,dx=0,Dx=10,dy=0,Dy=10](0,0)(-5,-5)(5,5)
\psplot[plotstyle=curve,arrows=<->]{-5}{2}{2 x exp -1 mul 2 add}
\end{pspicture}
&
\begin{pspicture}(-1.2,-1.2)(1.2,1.2)
\psset{xunit=0.2,yunit=0.2}
\psaxes[arrows=<->,dx=0,Dx=10,dy=0,Dy=10](0,0)(-5,-5)(5,5)
\psplot[plotstyle=curve,arrows=<->]{-5}{2}{2 x exp 2 add}
\end{pspicture}
\\\hline
$q<0$&


\begin{pspicture}(-1.2,-1.2)(1.2,1.2)
%\psgrid
\psset{xunit=0.2,yunit=0.2}
\psaxes[arrows=<->,dx=0,Dx=10,dy=0,Dy=10](0,0)(-5,-5)(5,5)
\psplot[plotstyle=curve,arrows=<->]{-5}{2}{2 x exp -1 mul 2 sub}
\end{pspicture}
&
\begin{pspicture}(-1.2,-1.2)(1.2,1.2)
%\psgrid
\psset{xunit=0.2,yunit=0.2}
\psaxes[arrows=<->,dx=0,Dx=10,dy=0,Dy=10](0,0)(-5,-5)(5,5)
\psplot[plotstyle=curve,arrows=<->]{-5}{2}{2 x exp 2 sub}
\end{pspicture}
\\\hline
\end{tabular}cos
\end{center}
\end{table}

\textbf{The effect of $q$}\newline

Jy sou gevind het dat die waarde van $q$ is called a vertical shift because all points are moved the same distance in the same direction (because it slides the entire graph up or down). 
\begin{itemize}
\item For $q>0$, then the graph is shifted vertically upwards by $q$ units. 
\item For $q<0$, then the graph is shifted vertically downwards by $q$ units.
\end{itemize}
The horizontal asymptote is shifted by $q$ units and is the line $y=q$. \vspace{8pt}\\


\textbf{The effect of $a$}\newline
Jy sou gevind het dat die waarde van $a$ bepaal die vorm van die grafiek. c
\begin{itemize}
 \item Indien $a>0$, dan sal die grafiek “Curve Upwards” – “CU”.
\item If $a<0$, dan sal die grafiek “Curves Downwards” – “CD” . It reflects the graph about the horizontal asymptote.
\end{itemize}

\subsection*{Discovering the characteristics}
algemene vorms en posisies van funksies van die vorm $y=ab^{x} + q$.
\subsection*{Definisieversameling en Waardeversameling}

Vir $y=ab^{x}+q$,is die funksie gedefinieer vir alle reële waardes van $x$.Dus, die definisieversameling is $\{x:x\in \mathbb{R}\}$.\par 
Die waardeversameling van $y=ab^{x}+q$ word bepaal deur die teken van $a$.\par 
Indien $a>0$ dan:\par
\begin{equation*}
\begin{array}{ccc}\hfill b^{x}& > & 0\hfill \\
 \hfill ab^{x}& > & 0\hfill \\ 
\hfill ab^{x}+q& > & q\hfill \\ 
\hfill f(x)& > & q\hfill 
\end{array}
\end{equation*}
Dus, as  $a>0$, dan is die waardeversameling  $\{f(x):f(x) > q\}$.\par 
Indien $a<0$ dan:\par 

\begin{equation*}
\begin{array}{ccc}\hfill b^{x}&< & 0\hfill \\
 \hfill ab^{x}& < & 0\hfill \\
\hfill ab^{x}+q& < & q\hfill \\
 \hfill f(x)& < & q\hfill 
\end{array}
\end{equation*}
Dus, as $a<0$, dan is die waardeversameling  $\{f(x):f(x) < q\}$.\par 

\begin{wex}{Domain and range of an exponential function}
{Find the domain and range of $g(x)=3.2^{x}+2$}
{
\westep{Find the domain}
die definisieversameling van $g(x)=3.2^{x}+2$ is $\{x:x\in \mathbb{R}\}$.
\westep{Find the range}
\begin{equation*}
\begin{array}{ccc}\hfill 2^{x}& > & 0\hfill \\
 \hfill 3.2^{x}& > & 0\hfill \\
 \hfill 3.2^{x}+2& > & 2\hfill 
\end{array}
\end{equation*}
Dus is die waardeversameling $\{g(x):g(x) > 2)\}$.\par 
}
 
\end{wex}




\subsection*{Afsnitte}
\textbf{The $y$-intercept}\\
Die $y$-afsnit word gegee deur $x=0$:
\begin{equation*}
\begin{array}{ccl}\hfill y& =& ab^{x}+q\hfill \\
 \hfill y}& =& ab^{0}+q\hfill \\
 & =& a(1)+q\hfill \\
 & =& a+q\hfill 
\end{array}
\end{equation*}

Byvoorbeeld, die $y$-afsnit van $g(x)=3.2^{x}+2$ word gegee deur 
 $x=0$ te stel, om dan te kry:\par 

\begin{equation*}
\begin{array}{ccl}\hfill y& =& 3.2^{x}+2\hfill \\
 \hfill y& =& 3.2^{0}+2\hfill \\
 & =& 3+2\hfill \\ & =& 5\hfill 
\end{array}
\end{equation*}
This gives the point $(0;5)$.\vspace{10pt}
\\
\textbf{The $x$-intercept}\\
Die $x$-afsnitte word bereken deur $y=0$ te stel. \\
Byvoorbeeld, die $x$-afsnit van $g(x)=3.2^{x}+2$ word gegee deur $y=0$ te stel:\par 
\begin{equation*}
\begin{array}{ccl}\hfill y& =& 3.2^{x}+2\hfill \\
 \hfill 0& =& 3.2^{x}+2\hfill \\
 \hfill -2& =& 3.2^{x}\hfill \\
 \hfill {2}^{x}& =& \dfrac{-2}{3}\hfill 
\end{array}
\end{equation*}
Dit het net ‘n rëele oplossing. Dus, het die grafiek van die vorm $g(x)$ geen $x$-afsnitte.\par 

\subsection*{Asimptote}

Daar is een asimptoot vir funksies van die vorm $y=ab^{x}+q$ Die horisontale asimptoot lê by $x=q$. 


\subsection*{Sketse van Grafieke van die vorm $f(x)=ab^{x}+q$}

Om grafieke te skets van funksies van die vorm, $f(x)=ab^{x}+q$, moet ons vier eienskappe bereken:\par 
\begin{enumerate}[noitemsep, label=\textbf{\arabic*}. ] 
\item shape
\item $y$-afsnit
\item $x$-afsnit
\item Asimptote
\end{enumerate}

\begin{wex}{Sketching an exponential function}
{ skets die grafiek van $g(x)=3(2^{x})+2$. Merk die afsnitte en asimptote.}
{
\westep{Examine the standard form of the equation}
From the equation we see that $a>1$, therefore the graph curves upwards. $q>0$ therefore the graph is shifted vertically upwards by $2$ units.

\westep{Calculate the intercepts}
Ons kry die $y$-afsnit, waar $x=0$:
\begin{equation*}
\begin{array}{ccl}\hfill y& =& 3.2^{x}+2\hfill \\
 \hfill y& =& 3.2^{0}+2\hfill \\
 & =& 3+2\hfill \\ & =& 5\hfill 
\end{array}
\end{equation*}
This gives the point $(0;5)$.\\

Ons kry die $x$-afsnit, waar $y=0$:
\begin{equation*}
\begin{array}{ccl}\hfill y& =& 3.2^{x}+2\hfill \\
 \hfill 0& =& 3.2^{x}+2\hfill \\
 \hfill -2& =& 3.2^{x}\hfill \\
 \hfill {2}^{x}& =& \dfrac{-2}{3}\hfill 
\end{array}
\end{equation*}
There is no real solution, en daar is geen $x$-afsnitte.

\westep{Determine the asymptote}
The horizontal asymptote is the line $y=2$.

\westep{Plot the points and sketch the graph}
\setcounter{subfigure}{0}
\begin{figure}[htbp]
\begin{center}
\begin{pspicture}(-5,-1)(5,6)
%\psgrid
\psset{yunit=0.75,xunit=0.75}
\psaxes[arrows=<->](0,0)(-5,-1)(5,7)
\psplot[plotstyle=curve,arrows=<->]{-5}{0.5}{2 x exp 3 mul 2 add}
\end{pspicture}
% \caption{Graph of $g(x)=3( 2^{x}) + 2$.}
% \label{fig:mf:g:exponentialsketchexample10}
\end{center}
\end{figure}    
\\
Domain: $\{x \in \mathBB{R}\}$\\
Range: $\{g(x): g(x) >2\}$\\

Note that there is no axis of symmetry for exponential functions.
} 
\end{wex}


 
\begin{wex}{Skets ’n Eksponsiële Grafiek}
{
Skets die grafieke van $y=-2.3^{x}+6$.}
{
\westep{Examine the standard form of the equation} 
From the equation we see that $a<0$ therefore the graph curves downwards. $q>0$ therefore the graph is shifted vertically upwards by $6$ units.
\westep{Calculate the intercepts}

Ons kry die $y$-afsnit waar $x=0$:
\begin{equation*}
\begin{array}{ccl}\hfill y& =& -2.3^{x}+6\hfill \\
 \hfill y& =& -2.3^{0}+6\hfill \\
 & =& 4\hfill \\

\end{array}
\end{equation*}
This gives the point $(0;4)$.\\

Ons kry die $x$-afsnit waar $y=0$:
\begin{equation*}
\begin{array}{ccl}\hfill y& =& -2.3^{x}+6\hfill \\
 \hfill 0& =& -2.3^{x}+6\hfill \\
 \hfill -6& =& -2.3^{x}\hfill \\
 \hfill {3}^{1}& =& {3}^{x}\hfill \\
\hfill \therefore x & =& 1 
\end{array}
\end{equation*}
This gives the point $(1; 0)$
\westep{Determine the asymptote} 
Funksies van hierdie vorm het een asimptoot. Dit lê by $y=q$ die asimptoot van die grafiek is by $y=5$.



\westep{Plot the points and sketch the graph} 
% Diagram needed here
\\
Die gebied is: $\{x \in \mathBB{R}\}$\\
die terrein is: $\{g(x): g(x) <6\}$\\

}
\end{wex}




\begin{exercises}{ }
 {
 
Skets die grafieke van $y=2^{x}$ en $y=(\frac{1}{2})^{x}$ op dieselfde assestelsel.
\begin{enumerate}[noitemsep, label=\textbf{\arabic*}. ] 
\item Is die $x$-as die asimptoot en/of simmetrie-as in albei grafieke? Verduidelik jou antwoord.
\item Watter grafiek word aangedui met die volgende vergelyking $y=2^{-x}$ ? Verduidelik jou antwoord.
\item Los die vergelyking $2^{x}=(\frac{1}{2})^{x}$ met behulp van ’n skets op en kontroleer jou antwoord deur middel van translasie.
\end{enumerate}
Die kurwe van die eksponensiele funksie $f$ in die meegaande diagram sny die y-as by die punt $A(0; 1)$ en $B(3; 9)$ is op $f$.
% THIS DIAGRAM NEEDS FIXING!
% \setcounter{subfigure}{0}
% \begin{center}
% \scalebox{1} % Change this value to rescale the drawing.
% {
% \begin{pspicture}(0,-3.19)(6.3184376,3.21)
% \psdots[dotsize=0.12](5.0,1.49)
% \psdots[dotsize= 0.12](1.0,-1.51)
% \psbezier[linewidth=0.04](0.0,-1.81)(5.0,-0.41)(5.5,2.59)(5.6,3.19)
% \rput(5.7725,1.7){$B(2;4)$}
% \rput(1.7076563,-1.8){$A(0;1)$}
% \rput(1.0,-2.51){\psaxes[linewidth=0.04,arrowsize=0.05291667cm 2.0,arrowlength= 1.4,arrowinset=0.4,ticksize=0.08cm,dx=2.0cm,dy=1.0cm]{<->}(0,0)(-1,0)(5,5)}
% \rput(5.994375,-2.8){$x$}
% \rput(0.564375,2.5){$f$}
% \end{pspicture} 
% }     
% \end{center}
\begin{enumerate}[noitemsep, label=\textbf{\arabic*}. ] 
\setcounter{enumi}{3}
\item  Bepaal die vergelyking van funksie $f$.
\item  Bepaal die vergelyking van $h$, die refleksie van die kurwe van $f$ in die $x$-as.
\item  Bepaal die waardeversameling van $h$.
\item Bepaal die vergelyking van  $g$, die refleksie van die kurwe van $f$ in die $y$-as.
\item Bepaal die vergelyking van $j$ indien $j$ is a vertical stretch of $f$ by $+2$ units.
\item Bepaal die vergelyking van $k$ indien $k$ is a vertical shift of $f$ by $-3$ units.
\end{enumerate}


\par \raisebox{-0.2em}{\includegraphics[height=1em]{../icons/www.pdf}} Kry die oplossing met die kortkodes:
\par \begin{tabular}[h]{cccccc}
(1.) lxk  &  (2.) lx0  & \end{tabular}
}
\end{exercises}

\section{Trigonometriese Funksies}
\nopagebreak
Hierdie afdeling beskryf die grafieke van trigonometriese funksies.\par 

\subsection{Grafiek van $y=sin~\theta $}
\subsection*{Plotting the graph}
\nopagebreak
Volgooi die volgende tabel en gebruik jou sakrekenaar om die waardes te bereken. Stip dan die waardes met $sin~\theta $ op die $y$-as en $\theta $op die $x$-as.Rond die antwoorde af tot $1$ desimale plek.\par 

\setlength\mytablespace{16\tabcolsep}
\addtolength\mytablespace{9\arrayrulewidth}
\setlength\mytablewidth{\linewidth}
\setlength\mytableroom{\mytablewidth}
\addtolength\mytableroom{-\mytablespace}
\setlength\myfixedwidth{0pt}
\setlength\mystarwidth{\mytableroom}
\addtolength\mystarwidth{-\myfixedwidth}
\divide\mystarwidth 80
% ----- Table with code
% \begin{table}[H]
% \\ '' '0'
\begin{center}
\label{m39414*id83562}
\noindent

\begin{tabular*}{\mytablewidth}{|p{10\mystarwidth}|p{10\mystarwidth}|p{10\mystarwidth}|p{10\mystarwidth}|p{10\mystarwidth}|p{10\mystarwidth}|p{10\mystarwidth}|p{10\mystarwidth}|}\hline

$\theta $  &
$0^{\circ }$ &
$30^{\circ }$ &
$60^{\circ }$ &
$90^{\circ }$ &
$120^{\circ }$ &
$150^{\circ }$ &
\\ \hline

$sin~\theta $ 
&
&
&
&
&
&
&
\\ \hline

$\theta $&
$180^{\circ }$ &
$210^{\circ }$ &
$240^{\circ }$ &
$270^{\circ }$ &
$300^{\circ }$ &
$330^{\circ }$ &
$360^{\circ }$% make-rowspan-placeholders
\\ \hline

$sin~\theta $
&
&
&
&
&
&
&
\\ \hline

\multicolumn{8}{|p{\dimexpr10\mystarwidth+10\mystarwidth+10\mystarwidth+10\mystarwidth+10\mystarwidth+10\mystarwidth+10\mystarwidth+10\mystarwidth+14\tabcolsep+7\arrayrulewidth\relax}|}{}

\\ \hline

\multicolumn{8}{|p{\dimexpr10\mystarwidth+10\mystarwidth+10\mystarwidth+10\mystarwidth+10\mystarwidth+10\mystarwidth+10\mystarwidth+10\mystarwidth+14\tabcolsep+7\arrayrulewidth\relax}|}{
\setcounter{subfigure}{0}
\label{m39414*id84030}
\begin{center}
\begin{pspicture}(0,-1)(4,1)
\psset{xunit=2}
%\psgrid[gridcolor=gray]
\psset{xunit=0.01111}
\psaxes[dx=30,Dx=30]{<->}(0,0)(0,-1.2)(370,1.2)
\end{pspicture}
\end{center}    
  }
\\ \hline
%--------------------------------------------------------------------
\end{tabular*}
\end{center}
\begin{center}{\small\bfseries Table 14.4}\end{center}
%\end{table}
\par
\label{m39414*id84056}Laat ons terugkyk na ons waardes vir $sin~\theta $\par 
\begin{table}[H]
% \begin{table}[H]
% \\ '' '0'
\begin{center}
\label{m39414*id84073}

\begin{tabular}{|l|l|l|l|l|l|l|}\hline
$\theta $&
${0}^{\circ }$&
${30}^{\circ }$&
${45}^{\circ }$&
${60}^{\circ }$&
${90}^{\circ }$&
${180}^{\circ }$
  % make-rowspan-placeholders
\\ \hline
%--------------------------------------------------------------------
$sin~\theta $&
$0$ &
$\frac{1}{2}$&
$\frac{1}{\sqrt{2}}$&
$\frac{\sqrt{3}}{2}$&
$1$ &
$0$% make-rowspan-placeholders
\\ \hline
%--------------------------------------------------------------------
\end{tabular}
\end{center}
\begin{center}{\small\bfseries Table 14.5}\end{center}
% \begin{caption}{\small\bfseries Table 14.5}\end{caption}
\end{table}
\par
Soos jy kan sien, die funksie $sin~\theta $ het ’n waarde van $0$ by $\theta ={0}^{\circ }$. Sy waarde neem egalig toe tot by $\theta ={90}^{\circ }$ wanneer sy waarde $1$.is. Ons weet ook dat dit later afneem na $0$ as $\theta ={180}^{\circ }$.Deur dit alles bymekaar te sit,
kan ons ’n idee kry van die volle omvang van die sinuskurwe. Die sinuskurwe word gewys in Figuur 14.23. Let
op die kurwe se vorm, waar elke kurwe die lengte het van ${360}^{\circ }$. Ons sê die grafiek het ’n periode van ${360}^{\circ }$.  Die
hoogte van die kurwe bo (of onder) die $x$-as word die kurwe se amplitude genoem. Dus is die amplitude van die
sinuskurwe is  $1$, and its minimum amplitude is $-1$.\par 
\setcounter{subfigure}{0}
\begin{figure}[h]
\begin{center}
\begin{pspicture}(-6,-2)(6,2)
\psaxes[Ox=0, Dx=180, dx=2]{<->}(0,0)(-4.5,-1.5)(4.5,1.5)
\psplot[xunit=0.0111,yunit=1.0, plotpoints=1000]{-360}{360}{x sin}
% \rput(5.3,-.3){$degrees$}
\psline[linestyle=dashed](0, 1)(4, 1)
\psline[linestyle=dashed](0, -1)(4, -1)
\end{pspicture}
\caption{Die grafiek van $sin~ \theta$.}
\label{trig:sin}
\end{center}
\end{figure}
      

\subsection*{Funksies in die vorm $y=a~sin~x+q$}
\nopagebreak
In die vergelyking, $y=a~sin~x+q$, $a$ en $q$  is konstantes en het verskillende invloede op die grafiek van die
funksie. Die algemene vorm van hierdie grafiek word gewys in Figuur 14.24 vir die funksie $f(\theta )=2sin~\theta +3$.\par 
\setcounter{subfigure}{0}
\begin{figure}[!ht]
\begin{center}
\begin{pspicture}(-4,-2)(4,6)
%\psgrid[gridcolor=gray]
\psset{yunit=1,xunit=0.01111}
\psaxes[dx=90,Dx=90]{<->}(0,0)(-360,-2)(360,6)
\psplot[plotstyle=curve,arrows=<->]{-360}{360}{x sin 2 mul 3 add}
% \rput(420,-0.3){$degrees$}
\end{pspicture}
\caption{Die grafiek van $f(\theta)=2 \sin \theta +3$}
\label{fig:mt:g:sin}
\end{center}
\end{figure}   

\subsection*{Investigating the effects of a and q}
\nopagebreak
\begin{enumerate}[noitemsep, label=\textbf{\arabic*}. ] 
\item Op dieselfde stel asse, trek die volgende grafieke:
\begin{enumerate}[noitemsep, label=\textbf{\alph*}. ] 
\item $a(\theta )=sin~\theta -2$
\item $b(\theta )=sin~\theta -1$
\item $c(\theta )=sin~\theta $
\item $d(\theta )=sin~\theta +1$
\item $e(\theta )=sin~\theta +2$
\end{enumerate}
Gebruik jou resultate om die invloed van $q$ af te lei.
\item Op dieselfde stel asse, trek die volgende grafieke:
\begin{enumerate}[noitemsep, label=\textbf{\alph*}. ] 
\item $f(\theta )=-2~sin~\theta $
\item $g(\theta )=-1~sin~\theta $
\item $h(\theta )=0~sin~\theta $
\item $j(\theta )=1~sin~\theta $
\item $k(\theta )=2~sin~\theta $\end{enumerate}
Gebruik jou resultate om die invloed van $a$ af te lei.
\end{enumerate}
Jy behoort te vind dat die waarde van $a$ die hoogte van die pieke van die grafiek beïnvloed. As die grootte van $a$ toeneem, word die pieke hoër. As dit afneem, word die pieke laer.\par 
$q$ is called the vertical shift. If $q=2$, then the whole sine graph shifts up $2$ units. If $q=-1$, the whole sine graph shifts down $1$ unit.\par 
Hierdie eienskappe word opgesom in Tabel 14.6.\par 
\begin{table}[htb]
\begin{center}
\caption{Tabel wat die algemene vorms en posisies van grafieke en funksies in die vorm $y=a \sin(x) + q$ opsom.}
\label{tab:mt:g:summarysin10}
\begin{tabular}{|c|c|c|}\hline
& $a>0$&$a<0$\\\hline
$q>0$&
\begin{pspicture}(-1.2,-0.6)(1.2,1.2)
\psset{yunit=0.5,xunit=0.0111}
\psaxes[arrows=<->,dx=0,Dx=720,dy=0,Dy=10,xunit=0.25](0,0)(-450,-1)(450,2)
\psplot[plotstyle=curve,arrows=<->,xunit=0.25]{-360}{360}{x sin 0.5 add}
% \rput(5.1,-.3){Degrees}
\end{pspicture}
&
\begin{pspicture}(-1.2,-0.6)(1.2,1.2)
\psset{yunit=0.5,xunit=0.0111}
\psaxes[arrows=<->,dx=0,Dx=720,dy=0,Dy=10,xunit=0.25](0,0)(-450,-1)(450,2)
\psplot[plotstyle=curve,arrows=<->,xunit=0.25]{-360}{360}{x sin neg 0.5 add}
% \rput(5.1,-.3){Degrees}
\end{pspicture}\\\hline
$q<0$&
\begin{pspicture}(-1.2,-1.2)(1.2,0.6)
%\psgrid
\psset{yunit=0.5,xunit=0.0111}
\psaxes[arrows=<->,dx=0,Dx=720,dy=0,Dy=10,xunit=0.25](0,0)(-450,-2)(450,1)
\psplot[plotstyle=curve,arrows=<->,xunit=0.25]{-360}{360}{x sin 0.5 sub}
% \rput(5.1,-.3){Degrees}
\end{pspicture}
&
\begin{pspicture}(-1.2,-1.2)(1.2,0.6)
%\psgrid
\psset{yunit=0.5,xunit=0.0111}
\psaxes[arrows=<->,dx=0,Dx=720,dy=0,Dy=10,xunit=0.25](0,0)(-450,-2)(450,1)
\psplot[plotstyle=curve,arrows=<->,xunit=0.25]{-360}{360}{x sin neg 0.5 sub}
% \rput(5.1,-.3){Degrees}
\end{pspicture}\\\hline
\end{tabular}
\end{center}
\end{table}
\par

\subsection*{Discovering the characteristics}
\subsection*{Gebied en Terrein}
\nopagebreak
Die $f(\theta )=a~sin~\theta +q$, is die gebied $\{\theta :\theta \in \mathbb{R}\}$  omdat daar geen waarde is van $\theta \in \mathbb{R}$ waarvoor $f(\theta )$ ongedefinieerd is nie..\par 
Die terrein van $f(\theta )=a~sin~\theta +q$ hang daarvan af of die waarde vir $a$ positief of negatief is. Ons sal die twee
gevalle afsonderlik oorweeg.\par 
As $a>0$ het ons:\par 
\nopagebreak\noindent{}
\begin{equation*}
\begin{array}{ccc}\hfill -1\le sin~\theta & \le & 1\hfill \\ \hfill -a\le a~sin~\theta & \le & a\hfill \\ \hfill -a+q\le a~sin~\theta +q& \le & a+q\hfill \\ \hfill -a+q\le f(\theta )& \le & a+q\hfill \end{array}
\end{equation*}
Dit vertel ons dat vir alle waardes van $\theta $, $f(\theta )$ altyd tussen $-a+q$ en $a+q$ is.\par
Daarom as $a>0$ is die terrein van  $f(\theta )=a~sin~\theta +q$ is $\{f(\theta ):f(\theta )\in [a+q,-a+q]\}$.
Insgelyks, daar kan getoon word dat as $a<0$, dan is die terrein van $f(\theta )=a~sin~\theta +q$ is $\{f(\theta ):f(\theta )\in [a+q,-a+q]\}$. Dit word as ’n oefening gelaat.\par 

\Tip{Die maklikste manier om die terrein te bepaal is om bloot vir die "bokant" en die "onderkant" van die grafiek te soek.}


\subsection*{Snypunte}
\nopagebreak
The $y$-snypunt, ${y}_{int}$, of $f(\theta )=a~sin~x+q$ is eenvoudig die waarde van$f(\theta )$ by $\theta ={0}^{\circ }$.\
\begin{equation*}
\begin{array}{ccc}\hfill {y}_{int}& =& f({0}^{\circ })\hfill \
\ & =& a~sin({0}^{\circ })+q\hfill \\
 & =& a(0)+q\hfill \\
 & =& q\hfill 
\end{array}
\end{equation*}



\subsection{Grafiek van $y=cos~\theta $}
\subsection*{Plotting the graph}
\nopagebreak
 Voltooi die volgende tabel, gebruik jou sakrekenaar om die waardes korrek tot $1$ desimale plek te bereken. Stip
dan die waardes met $cos~\theta $ op die $y$-as en $\theta $ op die $x$-as.\par 

% -------------------------
\setlength\mytablespace{16\tabcolsep}
\addtolength\mytablespace{9\arrayrulewidth}
\setlength\mytablewidth{\linewidth}
\setlength\mytableroom{\mytablewidth}
\addtolength\mytableroom{-\mytablespace}
\setlength\myfixedwidth{0pt}
\setlength\mystarwidth{\mytableroom}
\addtolength\mystarwidth{-\myfixedwidth}
\divide\mystarwidth 80subsub
% ----- Table with code
% \begin{table}[H]
% \\ '' '0'
\begin{center}
\label{m39414*id86399}
\noindent

\begin{tabular*}{\mytablewidth}{|p{10\mystarwidth}|p{10\mystarwidth}|p{10\mystarwidth}|p{10\mystarwidth}|p{10\mystarwidth}|p{10\mystarwidth}|p{10\mystarwidth}|p{10\mystarwidth}|}\hline
$\theta $     &
$0^{\circ }$ &
$30^{\circ }$ &
$60^{\circ }$ &
$90^{\circ }$ &
$120^{\circ }$ &
$150^{\circ }$ &
\\ \hline

$cos~\theta $  &
&
&
&
&
&
&
\\ \hline

$\theta $    &
$180^{\circ }$ &
$210^{\circ }$ &
$240^{\circ }$ &
$270^{\circ }$ &
$300^{\circ }$ &
$330^{\circ }$ &
$360^{\circ }$

\\ \hline

$cos~\theta $&
&
&
&
&
&
&

\\ \hline

\multicolumn{8}{|p{\dimexpr10\mystarwidth+10\mystarwidth+10\mystarwidth+10\mystarwidth+10\mystarwidth+10\mystarwidth+10\mystarwidth+10\mystarwidth+14\tabcolsep+7\arrayrulewidth\relax}|}{}

\\ \hline

\multicolumn{8}{|p{\dimexpr10\mystarwidth+10\mystarwidth+10\mystarwidth+10\mystarwidth+10\mystarwidth+10\mystarwidth+10\mystarwidth+10\mystarwidth+14\tabcolsep+7\arrayrulewidth\relax}|}{
\setcounter{subfigure}{0}
\begin{pspicture}(0,-1)(4,1)
\psset{xunit=2}
%\psgrid[gridcolor=gray]
\psset{xunit=0.01111}
\psaxes[dx=30,Dx=30]{<->}(0,0)(0,-1.2)(370,1.2)
\end{pspicture} 
  }

\\ \hline

\end{tabular*}
\end{center}
\begin{center}{\small\bfseries Table 14.7}\end{center}
%\end{table}
\par
\label{m39414*id86892}Laat ons terugkyk na ons waardes vir $cos~\theta $\par 
% \textbf{m39414*uid46}\par
\begin{table}[H]
% \begin{table}[H]
% \\ '' '0'
\begin{center}
\label{m39414*id86909}
\noindent

\begin{tabular}{|l|l|l|l|l|l|l|}\hline
$\theta $
${0}^{\circ }$&
${30}^{\circ }$&
${45}^{\circ }$&
${60}^{\circ }$&
${90}^{\circ }$&
${180}^{\circ }$&
  % make-rowspan-placeholders
\\ \hline
%--------------------------------------------------------------------
$cos~\theta $&
$1$ &
$\frac{\sqrt{3}}{2}$&
$\frac{1}{\sqrt{2}}$&
$\frac{1}{2}$&
$0$ &
$-1$
  % make-rowspan-placeholders
\\ \hline
%--------------------------------------------------------------------
\end{tabular}
\end{center}
\begin{center}{\small\bfseries Table 14.8}\end{center}
% \begin{caption}{\small\bfseries Table 14.8}\end{caption}
\end{table}
\par
As jy noukeurig kyk, sal jy oplet dat die cosinus van ’n hoek $\theta $  dieselfde is as die sinus van die hoek ${90}^{\circ }-\theta $. Neem byvoorbeeld,\par 
\nopagebreak\noindent{}
\begin{equation*}
cos~{60}^{\circ }=\frac{1}{2}=sin~{30}^{\circ }=sin~({90}^{\circ }-{60}^{\circ })
\end{equation*}
Dit wys ons dat ten einde ’n cosinusgrafiek te skep, al wat ons hoef te doen is om die sinusgrafiek ${90}^{\circ }$ na links te
skuif. die grafiek van $cos~\theta $ word gewys in Figuur 14.30. As die cosinusgrafiek eenvoudig ’n geskuifde sinusgrafiek
is, sal dit dieselfde periode en amplitude as die sinuskurwe hê.\par 
\setcounter{subfigure}{0}
\begin{figure}[h]
\begin{center}
\begin{pspicture}(-6,-2)(6,2)
\psaxes[Ox=0, Dx=180, dx=2]{<->}(0,0)(-4.5,-1.5)(4.5,1.5)
%\psline[]{<-}(0,1.2)(1,1.2)
\psplot[xunit=0.0111,yunit=1.0, plotpoints=1000]{-360}{360}{x cos}
% \rput(5.1,-.3){Degrees}
\end{pspicture}
\caption{The graph of $\cos \theta$.}
\label{trig:cos}
\end{center}
\end{figure}      

\subsection*{Funksies in die vorm  $y=a~cos~x+q$}
\nopagebreak
In die vergelyking, $y=a~cos~x+q$, $a$ en $q$ is konstantes en het verskillende invloede op die grafiek van die
funksie. Die algemene vorm van die grafieke van hierdie soort funksies word getoon in Figuur 14.31 vir die
funksie $f(\theta )=2~cos~\theta +3$.\par 
\setcounter{subfigure}{0}
\begin{figure}[!ht]
\begin{center}
\begin{pspicture}(-4,-2)(4,6)
%\psgrid[gridcolor=gray]
\psset{yunit=1,xunit=0.01111}
\psaxes[dx=90,Dx=90]{<->}(0,0)(-360,-2)(360,6)
\psplot[plotstyle=curve,arrows=<->]{-360}{360}{x cos 2 mul 3 add}
% \rput(5.1,-.3){Degrees}
\end{pspicture}
\caption{Graph of $f(\theta)=2 ~\cos ~\theta +3$}
\label{fig:mt:g:cos}
\end{center}
\end{figure}      

\subsection*{Investigating the effects of a and q}
\nopagebreak
\begin{enumerate}[noitemsep, label=\textbf{\arabic*}. ] 
\item Op dieselfde stel asse, trek die volgende grafieke:
\begin{enumerate}[noitemsep, label=\textbf{\alph*}. ] 
\item $a(\theta )=cos~\theta -2$
\item $b(\theta )=cos~\theta -1$
\item $c(\theta )=cos~\theta $
\item $d(\theta )=cos~\theta +1$
\item $e(\theta )=cos~\theta +2$\end{enumerate}
Gebruik jou resultate om die invloed van $q$ af te lei.
\item Op dieselfde stel asse, trek die volgende grafieke:
\begin{enumerate}[noitemsep, label=\textbf{\alph*}. ] 
\item $f(\theta )=-2~cos~\theta $
\item $g(\theta )=-1~cos~\theta $
\item $h(\theta )=0~cos~\theta $
\item $j(\theta )=1~cos~\theta $
\item $k(\theta )=2~cos~\theta $\end{enumerate}
Gebruik jou resultate om die invloed van $a$ af te lei.
\end{enumerate}
Ons vind dat die waarde van $a$ die amplitude van die cosinusgrafiek op dieselfde manier beïnvloed as wat dit vir
die sinusgrafiek gedoen het.\par 
Verandering in die waarde van $q$ sal die die cosinusgrafiek op dieselfde manier skuif as wat dit vir die sinusgrafiek
gedoen het.\par 
Die verskillende eienskappe word opgesom in Tabel 14.9.\par 
\begin{table}[htb]
\begin{center}
\caption{Tabel wat die algemene vorms en posisies van grafieke en funksies in die vorm $y=a~ \cos~x + q$ opsom.}
\label{tab:mt:g:summarycos10}
\begin{tabular}{|c||c|c|}\hline
& $a>0$&$a<0$\\\hline\hline
$q>0$&
\begin{pspicture}(-1.2,-0.6)(1.2,1.2)
\psset{yunit=0.5,xunit=0.0111}
\psaxes[arrows=<->,dx=0,Dx=720,dy=0,Dy=10,xunit=0.25](0,0)(-450,-1)(450,2)
\psplot[plotstyle=curve,arrows=<->,xunit=0.25]{-360}{360}{x cos 0.5 add}
\end{pspicture}
&
\begin{pspicture}(-1.2,-0.6)(1.2,1.2)
\psset{yunit=0.5,xunit=0.0111}
\psaxes[arrows=<->,dx=0,Dx=720,dy=0,Dy=10,xunit=0.25](0,0)(-450,-1)(450,2)
\psplot[plotstyle=curve,arrows=<->,xunit=0.25]{-360}{360}{x cos neg 0.5 add}
\end{pspicture}\\\hline
$q<0$&
\begin{pspicture}(-1.2,-1.2)(1.2,0.6)
%\psgrid
\psset{yunit=0.5,xunit=0.0111}
\psaxes[arrows=<->,dx=0,Dx=720,dy=0,Dy=10,xunit=0.25](0,0)(-450,-2)(450,1)
\psplot[plotstyle=curve,arrows=<->,xunit=0.25]{-360}{360}{x cos 0.5 sub}
\end{pspicture}
&
\begin{pspicture}(-1.2,-1.2)(1.2,0.6)
%\psgrid
\psset{yunit=0.5,xunit=0.0111}
\psaxes[arrows=<->,dx=0,Dx=720,dy=0,Dy=10,xunit=0.25](0,0)(-450,-2)(450,1)
\psplot[plotstyle=curve,arrows=<->,xunit=0.25]{-360}{360}{x cos neg 0.5 sub}
\end{pspicture}\\\hline
\end{tabular}
\end{center}
\end{table}
\par

\subsection*{Discovering the characteristics}
\subsection*{Gebied en Terrein}
\nopagebreak
Vir $f(\theta )=a~cos(\theta )+q$, is die gebied $\{\theta :\theta \in \mathbb{R}\}$ want daar is geen waarde van $\theta \in \mathbb{R}$ waarvoor $f(\theta )$ ongedefinieërd is nie.\par 
Dit is maklik om te sien dat die terrein van $f(\theta )$ dieselfde sal wees as die terrein van $a~sin(\theta )+q$. Dit is omdat
die maksimum en minimumwaardes van $a~cos(\theta )+q$ dieselfde is as die maksimum en minimumwaardes van $a~sin(\theta )+q$.\par 

\subsection*{Snypunte}
\nopagebreak
Die $y$-afsnit van $f(\theta )=a~cos~x+q$ word bereken op dieselfde wyse as vir sinus.\par 
\nopagebreak\noindent{}
\begin{equation*}
\begin{array}{ccc}\hfill {y}_{int}& =& f({0}^{\circ })\hfill \\
 & =& a~cos({0}^{\circ })+q\hfill \\
 & =& a(1)+q\hfill \\
 & =& a+q\hfill 
\end{array}
\end{equation*}

\subsection*{Vergelyking van die Grafieke van $y=sin~\theta $ en $y=cos~\theta $}
\nopagebreak
\setcounter{subfigure}{0}
\begin{figure}[h]
\begin{center}
\begin{pspicture}(-6,-2)(6,2)
\psaxes[Ox=0, Dx=180, dx=2]{<->}(0,0)(-4.5,-1.5)(4.5,1.5)
\psline[]{<-}(0,1.2)(1,1.2)
\psplot[xunit=0.0111, , yunit=1.0, plotpoints=1000, linestyle=dashed]{-360}{360}{x sin}
\psplot[xunit=0.0111,yunit=1.0, plotpoints=1000]{-360}{360}{x cos}
% \rput(5.1,-.3){Degrees}
\rput(1.8,1.2){$90^\circ$ shift}
\end{pspicture}
\caption{Die grafiek van $\cos \theta$ (soliede lyn) en die grafiek van $\sin \theta$ (stippellyn).}
\end{center}
\end{figure}    
Let daarop dat die twee grafieke baie eenders lyk. Beide ossilleer op en af rondom die $x$-as soos wat jy beweeg
langs die as. Die afstande tussen die pieke van die twee grafieke is dieselfde en is konstant vir elke grafiek. Die
hoogte van elke piek en die diepte van elke trog is dieselfde.\par 
Die enigste verskil is dat die $sin$  grafiek skuif ’n bietjie na regs ten opsigte van die $cos$ grafiek, met $90^{\circ }$. Dit beteken
dat as ons die hele cos grafiek $cos$ grafiek $90 ^{\circ }$ na regs skuif, sal dit perfek oorvleul met die $sin$ grafiek. Jy kan ook die $sin$ grafiek $90 ^{\circ }$na links skuif en dan sal dit perfek oorvleul met die $cos$  grafiek. Dit beteken dat:\par 
\nopagebreak\noindent{}
\begin{equation*}
\begin{array}{ccc}\hfill sin~\theta & =& cos~(\theta -90)(\mbox{skuif die } cos{ grafiek na die regterkant})\hfill \\
 & \mathbf{a}\mbox{nd}& \\
 \hfill cos~\theta & =& sin(\theta +90)(\mbox{skuif die}sin\mbox{ grafiek na die linkerkant})\hfill 
\end{array}
\end{equation*}

\subsection{Grafiek van $y=tan~\theta $ }
\subsection*{Plotting the graph}
\nopagebreak
Voltooi die volgende tabel, gebruik jou sakrekenaar en bereken die waardes korrek tot $1$ desimale plek. Stip dan
die waardes met $tan~\theta $ op die $y$-as en $\theta $ op die $x$-as.\par 

% -------------------------
\setlength\mytablespace{16\tabcolsep}
\addtolength\mytablespace{9\arrayrulewidth}
\setlength\mytablewidth{\linewidth}
\setlength\mytableroom{\mytablewidth}
\addtolength\mytableroom{-\mytablespace}
\setlength\myfixedwidth{0pt}
\setlength\mystarwidth{\mytableroom}
\addtolength\mystarwidth{-\myfixedwidth}
\divide\mystarwidth 80

\begin{center}
\label{m39414*id89083}
\noindent

\begin{tabular*}{\mytablewidth}{|p{10\mystarwidth}|p{10\mystarwidth}|p{10\mystarwidth}|p{10\mystarwidth}|p{10\mystarwidth}|p{10\mystarwidth}|p{10\mystarwidth}|p{10\mystarwidth}|}\hline

$\theta $ &
$0^{\circ }$ &
$30^{\circ }$ &
$60^{\circ }$ &
$90^{\circ }$ &
$120^{\circ }$ &
$150^{\circ }$ &

\\ \hline

$tan~\theta $ &
&
&
&
&
&
&

\\ \hline

$\theta $ &
$180^{\circ }$ &
$210^{\circ }$ &
$240^{\circ }$ &
$270^{\circ }$ &
$300^{\circ }$ &
$330^{\circ }$ &
$360^{\circ }$
\\ \hline

$tan~\theta $ &
&
&
&
&
&
&

\\ \hline
%--------------------------------------------------------------------

\multicolumn{8}{|p{\dimexpr10\mystarwidth+10\mystarwidth+10\mystarwidth+10\mystarwidth+10\mystarwidth+10\mystarwidth+10\mystarwidth+10\mystarwidth+14\tabcolsep+7\arrayrulewidth\relax}|}{}

\\ \hline

\multicolumn{8}{|p{\dimexpr10\mystarwidth+10\mystarwidth+10\mystarwidth+10\mystarwidth+10\mystarwidth+10\mystarwidth+10\mystarwidth+10\mystarwidth+14\tabcolsep+7\arrayrulewidth\relax}|}{
\setcounter{subfigure}{0}
\begin{pspicture}(0,-1)(4,1)
\psset{xunit=2}
%\psgrid[gridcolor=gray]
\psset{xunit=0.01111}
\psaxes[dx=30,Dx=30]{<->}(0,0)(0,-1.2)(370,1.2)
\end{pspicture}   
  }

\\ \hline
%--------------------------------------------------------------------
\end{tabular*}
\end{center}
\begin{center}{\small\bfseries Table 14.10}\end{center}
%\end{table}
\par
\label{m39414*id89576}Kom ons kyk weer na ons waardes vir $tan~\theta $\par 
% \textbf{m39414*uid65}\par
\begin{table}[H]
% \begin{table}[H]
% \\ '' '0'
\begin{center}
\label{m39414*id89593}
\noindent

\begin{tabular}[t]{|l|l|l|l|l|l|l|}\hline
$\theta $&
${0}^{\circ }$&
${30}^{\circ }$&
${45}^{\circ }$&
${60}^{\circ }$&
${90}^{\circ }$&
${180}^{\circ }$
\\ \hline
%--------------------------------------------------------------------
$tan~\theta $&
$0$ &
$\frac{1}{\sqrt{3}}$&
$1$ &
$\sqrt{3}$&
$\infty $&
$0$% make-rowspan-placeholders
\\ \hline
%--------------------------------------------------------------------
\end{tabular}
\end{center}
\begin{center}{\small\bfseries Table 14.11}\end{center}
% \begin{caption}{\small\bfseries Table 14.11}\end{caption}
\end{table}
\par
Nou dat ons die grafieke het vir $sin~\theta $ en $cos~\theta $, is daar ’n maklike manier om die tan-grafiek te visualiseer. Kom ons
kyk weer na ons definisies van $sin~\theta $ en $cos~\theta $ vir ’n reghoekige driehoek.\par 
\nopagebreak\noindent{}
\begin{equation*}
\frac{sin~\theta }{cos~\theta }=\frac{\frac{\mbox{opposite}}{\mbox{hypotenuse}}}{\frac{\mbox{adjacent}}{\mbox{hypotenuse}}}=\frac{\mbox{opposite}}{\mbox{adjacent}}=tan~\theta 
\end{equation*}
Dit is die eerste van ’n stel belangrike verbande wat ons trigonometriese identiteite noem. ’n Identiteit is waar vir
enige waarde van die onbekende(s) wat daarin ingestel word. In hierdie geval het ons aangetoon dat\par 
\nopagebreak\noindent{}
\begin{equation*}
tan~\theta =\frac{sin~\theta }{cos~\theta }
\end{equation*}
vir enige waarde van $\theta $.\par 
Dus weet ons dat vir die waardes van $\theta $  waarvoor $sin~\theta =0$, moet ook  $tan~\theta =0$.Soortgelyk, as $cos~\theta =0$ s die
waarde van $tan~\theta $  ongedefiniëerd omdat ons nie mag deel met $0$ nie. Die grafiek word getoon in Figuur 14.38. Die
vertikale stippellyne is die waardes van $\theta $  waarvoor $tan~\theta $ nie gedefiniëerd is nie.\par 
\setcounter{subfigure}{0}
\\begin{figure}[h]
\begin{center}
\begin{pspicture}(-6,-3)(6,3)
\psaxes[Dx=180, dx=2, Dy=2, dy=1]{<->}(0,0)(-4.5,-3)(4.5,3)
\psline[linestyle=dashed](-1,-2.5)(-1,2.5)
\psline[linestyle=dashed](1,-2.5)(1,2.5)
\psline[linestyle=dashed](-3,-2.5)(-3,2.5)
\psline[linestyle=dashed](3,-2.5)(3,2.5)
\psplot[xunit=0.0111,yunit=0.5, plotpoints=500, arrows=<->]{-80}{80}{x sin x cos div}
\psplot[xunit=0.0111,yunit=0.5,plotpoints=500, arrows=<->]{-260}{-100}{x sin x cos div}
\psplot[xunit=0.0111,yunit=0.5,plotpoints=500, arrows=<->]{260}{100}{x sin x cos div}
\psplot[xunit=0.0111,yunit=0.5,plotpoints=500, arrows=<->]{-100}{-260}{x sin x cos div}
\psplot[xunit=0.0111,yunit=0.5,plotpoints=500, arrows=<-]{-280}{-360}{x sin x cos div}
\psplot[xunit=0.0111,yunit=0.5,plotpoints=500, arrows=<-]{280}{360}{x sin x cos div}
% \rput(5.1,-.3){$\theta$\ Degrees}
\rput(.4,3.3){$f(\theta$)}
\end{pspicture}
% \caption{Die grafiek van $\tan \theta$.}
\label{trig:tan}
\end{center}
\end{figure}       

\subsection*{Funksies van die vorm $y=a ~tan~x+q$}
\nopagebreak
Die figuur hieronder is ’n voorbeeld van ’n funksie van die vorm $y=a ~tan(x)+q$.\par 
\setcounter{subfigure}{0}
\begin{figure}[!ht]
\begin{center}
\begin{pspicture}(-5,-3)(5,3.2)
%\psgrid
\psset{yunit=0.25}
\psaxes[Dx=180, dx=2, Dy=5, dy=5]{<->}(0,0)(-4.5,-12)(4.5,12)
\psline[linestyle=dashed](-1,-12.5)(-1,12.5)
\psline[linestyle=dashed](1,-12.5)(1,12.5)
\psline[linestyle=dashed](-3,-12.5)(-3,12.5)
\psline[linestyle=dashed](3,-12.5)(3,12.5)
\psplot[xunit=0.0111, plotpoints=500, arrows=<->]{-80}{80}{x sin x cos div 2 mul 1 add}
\psplot[xunit=0.0111,plotpoints=500, arrows=<->]{-260}{-100}{x sin x cos div 2 mul 1 add}
\psplot[xunit=0.0111,plotpoints=500, arrows=<->]{260}{100}{x sin x cos div 2 mul 1 add}
\psplot[xunit=0.0111,plotpoints=500, arrows=<->]{-100}{-260}{x sin x cos div 2 mul 1 add}
\psplot[xunit=0.0111,plotpoints=500, arrows=<-]{-280}{-360}{x sin x cos div 2 mul 1 add}
\psplot[xunit=0.0111,plotpoints=500, arrows=<-]{280}{360}{x sin x cos div 2 mul 1 add}
% \rput(5.3,-.6){$\theta$\ Degrees}
\rput(.4,9.3){$f(\theta$)}
\end{pspicture}
\caption{The graph of $2\tan \theta + 1$.}
\label{trig:tan2}
\end{center}
\end{figure}
       

\subsection*{Investigating the effects of a and q}
\nopagebreak
\begin{enumerate}[noitemsep, label=\textbf{\arabic*}. ] 
\item Op dieselfde assestelsel, trek die volgende grafieke::
\begin{enumerate}[noitemsep, label=\textbf{\alph*}. ] 
\item $a(\theta )=tan~\theta -2$
\item $b(\theta )=tan~\theta -1$
\item $c(\theta )=tan~\theta $
\item $d(\theta )=tan~\theta +1$
\item $e(\theta )=tan~\theta +2$\end{enumerate}
Gebruik jou resultate om die invloed van $q$ af te lei.
\item Op dieselfde assestelsel, trek die volgende grafieke:
\begin{enumerate}[noitemsep, label=\textbf{\alph*}. ] 
\item $f(\theta )=-2~tan~\theta $
\item $g(\theta )=-1~tan~\theta $
\item $h(\theta )=0~tan~\theta $
\item $j(\theta )=1~tan~\theta $
\item $k(\theta )=2~tan~\theta $\end{enumerate}
Gebruik jou resultate om die invloed van $a$ af te lei.
\end{enumerate}
Ons vind dat die waarde van $a$ die steilheid van die bene van die grafiek beinvloed. Hoe groter die absolute waarde vanf $a$, hoe vinniger nader die bene die waardes van hulle asimptote, die waardes waar hulle nie gedefinieërd is nie. Negatiewe $a$ waardes keer die rigting waarin die bene van die grafiek loop, om. Ons vind
verder dat die waarde van $q$ beïnvloed die vertikale verskuiwing net soos by $sin~\theta $ en $cos~\theta $.
Hierdie verskillende eienskappe word opgesom in Tabel 14.12.\par 
\begin{table}[htb]
\begin{center}
\caption{Tabel van die algemene vorms en posisies van grafieke en funksies van die vorm $y=a~ \tan~x + q$.}
\label{tab:mt:g:summarytan10}
\begin{tabular}{|c||c|c|}\hline
& $a>0$&$a<0$\\\hline\hline
$q>0$&
\begin{pspicture}(-1.2,-0.6)(1.2,0.8)
%\psgrid[gridcolor=gray]
\psset{yunit=0.1,xunit=0.0111}
\psaxes[arrows=<->,dx=0,Dx=720,dy=0,Dy=10,xunit=0.25](0,0)(-450,-6)(450,7)
\psplot[plotstyle=curve,arrows=<->,xunit=0.25]{-81.5}{78}{x sin x cos div 1.5 add}cos
\end{pspicture}
&
\begin{pspicture}(-1.2,-0.6)(1.2,0.8)
%\psgrid[gridcolor=gray]
\psset{yunit=0.1,xunit=0.0111}
\psaxes[arrows=<->,dx=0,Dx=720,dy=0,Dy=10,xunit=0.25](0,0)(-450,-6)(450,7)
\psplot[plotstyle=curve,arrows=<->,xunit=0.25]{-78}{82.5}{x sin x cos div neg 1.5 add}
\end{pspicture}\\\hline
$q<0$&
\begin{pspicture}(-1.2,-0.8)(1.2,0.8)
%\psgrid[gridcolor=gray]
\psset{yunit=0.1,xunit=0.0111}
\psaxes[arrows=<->,dx=0,Dx=720,dy=0,Dy=10,xunit=0.25](0,0)(-450,-7)(450,6)
\psplot[plotstyle=curve,arrows=<->,xunit=0.25]{-80}{80}{x sin x cos div 1.5 sub}
\end{pspicture}
&
\begin{pspicture}(-1.2,-0.8)(1.2,0.8)
%\psgrid
\psset{yunit=0.1,xunit=0.0111}
\psaxes[arrows=<->,dx=0,Dx=720,dy=0,Dy=10,xunit=0.25](0,0)(-450,-7)(450,6)
\psplot[plotstyle=curve,arrows=<->,xunit=0.25]{-80}{80}{x sin x cos div neg 1.5 sub}
\end{pspicture}\\\hline
\end{tabular}
\end{center}
\end{table}
\par

\subsection*{Discovering the characteristics}
\subsection*{Gebied en Terrein}
\nopagebreak
Die gebied van $f(\theta )=a~tan(\theta )+q$ is al die waardes van $\theta $ sodat $cos~\theta $ nie gelyk is aan $0$nie. Ons het reeds gesien dat as $cos~\theta =0$, $tan~\theta =\frac{sin~\theta }{cos~\theta }$ ongedefinieerd is, want ons het deling deur nul. Ons weet dat $cos~\theta =0$ for all $\theta ={90}^{\circ }+{180}^{\circ }n$, waar $n$ waar n ’n heelgetal is. Dus die gebied van $f(\theta )=a~tan(\theta )+q$ is alle waardes van $\theta $, behalwe die waardes $\theta ={90}^{\circ }+{180}^{\circ }n$.\par 
Die terrein van $f(\theta )=a~tan~\theta +q$ is $\{f(\theta ):f(\theta )\in (-\infty ,\infty )\}$.\par 

\subsection*{Snypunte}
\nopagebreak
Die $y$-snypunt, ${y}_{int}$, van $f(\theta )=a~tan~x+q$ is slegs die waarde van $f(\theta )$ by $\theta ={0}^{\circ }$.\par 
\nopagebreak\noindent{}
\begin{equation*}
\begin{array}{ccc}\hfill {y}_{int}& =& f({0}^{\circ })\hfill \\
 & =& a~tan({0}^{\circ })+q\hfill \\
 & =& a(0)+q\hfill \\
 & =& q\hfill 
\end{array}
\end{equation*}

\subsection*{Asimptote}
\nopagebreak
Soos $\theta $ geleidelik naderkom aan ${90}^{\circ }$,sal $tan~\theta $ nader kom aan oneindig. Maar omdat $\theta $ ongedefinieërd is by ${90}^{\circ }$,kan $\theta $ slegs al nader kom aan ${90}^{\circ }$, maar nooit daarby uitkom nie. So, die $tan~\theta $ grafiek kom nader en nader aan
die lyn  $\theta ={90}^{\circ }$,sonder om dit ooit te ontmoet. Dus die lyn $\theta ={90}^{\circ }$ is ’n asimptoot van $tan~\theta $. $tan~\theta $ het ook
asimptote by $\theta ={90}^{\circ }+{180}^{\circ }n$, waar $n$ ’n heelgetal is.\par 

\begin{exercises}{ Trigonometriese Funksies}
 {

\begin{enumerate}[noitemsep, label=\textbf{\arabic*}. ] 
\item Deur you kennis van die invloed van $a$ en $q$, te gebruik, skets elk van die volgende grafieke, sonder om ’n
tabel van waardes te gebruik, vir $\theta \in [{0}^{\circ };{360}^{\circ }]$
\begin{enumerate}[noitemsep, label=\textbf{\alph*}. ] 
\item $y=2~sin~\theta $
\item $y=-4~cos~\theta $
\item $y=-2~cos~\theta +1$
\item $y=sin~\theta -3$
\item $y=tan~\theta -2$\item $y=2~cos~\theta -1$
\end{enumerate}
 \item Gee die vergelykings van elk van die volgende grafieke:
\setcounter{subfigure}{0}
% \begin{minipage}{0.6\textwidth}
\begin{center}
\begin{pspicture}(-2.5,-2)(5,2)
\psset{yunit=0.25}
\psaxes[Dx=180, dx=2, Dy=2, dy=4]{<->}(0,0)(-2,-5.1)(4.5,5.1)
\psplot[xunit=0.0111, plotpoints=500, arrows=<->]{-90}{360}{x cos -4 mul }
\uput[d](4.7,0){$x$}
\uput[r](0,5.1){$y$}
% \rput(-2.5,2){a)}
\end{pspicture}

\begin{pspicture}(-0.3,-2)(5,2)
%\psgrid
\psset{yunit=0.25}
\psaxes[Dx=90, dx=1, Dy=2, dy=4]{<->}(0,0)(0,-5.1)(4.5,5.1)
\psplot[xunit=0.0111, plotpoints=500, arrows=->]{0}{360}{x sin 1 add 2 mul}
\uput[d](4.7,0){$x$}
\uput[r](0,5.1){$y$}
% \rput(-1.2,3){b)}
%\psplot[xunit=0.0111,plotpoints=500, arrows=<->]{-260}{-100}{x sin x cos div 2 mul 1 add}
%\psplot[xunit=0.0111,plotpoints=500, arrows=<->]{260}{100}{x sin x cos div 2 mul 
\end{pspicture}
\end{center}
% \end{minipage}
% \begin{minipage}{0.35\textwidth}
\begin{pspicture}(-2.2,-3)(2.2,3.2)
%\psgrid
\psset{yunit=0.2}
\psaxes[Dx=90, dx=1, Dy=5, dy=5]{<->}(0,0)(-2,-12)(2,12)
\psline[linestyle=dashed](-1,0)(-1,12.5)
\psline[linestyle=dashed](1,-12.5)(1,-3)
\psline[linestyle=dashed](-1,-12.5)(-1,-3)
\psline[linestyle=dashed](1,0)(1,12.5)
% \rput(-2,8){c)}
\psplot[xunit=0.0111, plotpoints=500, arrows=<->]{-75}{83}{x sin x cos div -2 mul 5 add}
\end{pspicture}
% \end{minipage}      \par 

\par \raisebox{-0.2em}{\includegraphics[height=1em]{../icons/www.pdf}} Kry die oplossing met die kortkodes:
\par \begin{tabular}[h]{cccccc}
(1.) la8  &  (2.) la0  & \end{tabular}
}
\end{exercises}

\section{Opsomming}


\summary

\begin{itemize}[noitemsep]
\item jy behoort die volgende kenmerke van funksies te ken: 
\begin{itemize}[noitemsep]
\item Die gegewe of gekose $x$-waarde is bekend as die onafhanklike veranderlike want die waarde van x kan
vrylik gekies word. Die berekende $y$-waarde staan bekend as die afhanklike veranderlike aangesien
die waarde van y afhang van die gekose waarde van $x$
\item Die gebied van ’n relasie is die versameling van al die $x$ waardes waarvoor daar ten minste een $y$ waarde bestaan volgens die funksievoorskrif. Die terrein is die versameling van al die $y$ values,
 which verkry kan word deur ten minste een van die $x$ waardes te gebruik.
\item Die afsnit is die punt waar die grafiek ’n as sny. Die $x$-afsnit(te) is die punt(e) waar die grafiek die $x$-as
sny en die $y$-afsnit(te) is die punt(e) waar die grafiek die $y$-as sny. 
\item Slegs vir grafieke van funksies met ’n hoogste mag van groter as 1. Daar is twee tipes draaipunte: ’n
minimum draaipunt en ’n maksimum draaipunt. ’n Minimum draaipunt is ’n punt op die grafiek waar
die grafiek ophou afneem in waarde en begin toeneem in waarde. ’n Maksimum draaipunt is ’n punt
op die grafiek waar die grafiek ophou toeneem in waarde en begin afneem in waarde. 
\item ’n Asimptoot is ’n reguitlyn of kurwe wat die grafiek van ’n funksie sal nader, maar nooit raak nie.
\item ’n Lyn ten opsigte waarvan die grafiek simmetries is.
\item Die interval waar die grafiek toeneem of afneem
\item n Grafiek is kontinu as daar geen onderbreking in die grafiek is nie. 
\end{itemize}
\item 
Versamelingnotasie: ’n versameling van sekere x-waardes het die volgende notasie: \{$x$ : voorwaardes,
meer voorwaardes\}
\item 

Interval notasie: hier skryf ons ’n interval in die vorm ’laer hakie, laer getal, kommapunt, hoër getal, hoër
hakie’
\item  Jy moet die volgende funksies en hulle eienskappe ken:
    \begin{itemize}[noitemsep]
    \item Funksies van die vorm $y=ax+q$. Dit is reguitlyne.
    \item Funksies van die vorm $y=a{x}^{2}+q$ Dit staan bekend as paraboliese funksies of parabole.
    \item Funksies van die vorm $y=\frac{a}{x}+q$. Dit staan bekend as hiperboliese funksies of hiperbole.
    \item Funksies van die vorm $y=a{b}^{(x)}+q$. Hulle staan bekend as eksponensiële funksies.
    \end{itemize}
\end{itemize}

\begin{eocexercises}{}
\nopagebreak
\begin{enumerate}[noitemsep, label=\textbf{\arabic*}. ] 
\item Skets die volgende straight lines: 
    \begin{enumerate}[noitemsep, label=\textbf{\alph*}. ] 
    \item $y=2x+4$ 
    \item $y-x=0$ 
    \item $y=-\frac{1}{2}x+2$
    \end{enumerate}
\item Skets die volgende funksies: 
    \begin{enumerate}[noitemsep, label=\textbf{\alph*}. ] 
    \item $y={x}^{2}+3$ 
    \item $y=\frac{1}{2}{x}^{2}+4$
    \item $y=2{x}^{2}-4$
    \end{enumerate}
\item Skets die volgende funksies and identify the asymptotes: 
    \begin{enumerate}[noitemsep, label=\textbf{\alph*}. ] 
    \item $y={3}^{x}+2$ 
    \item $y=-4.{2}^{x}+1$ 
    \item $y=2.{3}^{x}-2$ 
    \end{enumerate}
\item Skets die volgende funksies and identify the asymptotes: 
    \begin{enumerate}[noitemsep, label=\textbf{\alph*}. ] 
    \item $y=\frac{3}{x}+4$ 
    \item $y=\frac{1}{x}$ 
    \item $y=\frac{2}{x}-2$ 
    \end{enumerate}
\item Determine whether the following statements are true or false. If the statement is false, give reasons why:
    \begin{enumerate}[noitemsep, label=\textbf{\alph*}. ] 
    \item  Die gegewe of gekose y-waarde staan bekend as die afhanklike veranderlike.
    \item Die afsnit is die punt waar die grafiek ’n as sny.
    \item There are two types of turning points -- minimal and maximal.
    \item A graph is said to be congruent if there are no breaks in the graph.
    \item Funksies van die vorm $y=ax+q$ is reguitlyne.
    \item Funksies van die vorm $y=\frac{a}{x}+q$ is eksponensiële funksies.
    \item  An asymptote is a straight or curved line which a graph will intersect once.
    \item Gegee die funksie in die vorm $y=ax+q$ , to find the y-intersect put $x=0$ and solve for $y$.
    \item The graph of a straight line always has a turning point.
    \end{enumerate}
\item Gegee die funksies $f(x)=-2{x}^{2}-18$ en $g(x)=-2x+6$
    \begin{enumerate}[noitemsep, label=\textbf{\alph*}. ] 
    \item Skets $f$ en $g$ op dieselfde assestelsel.
    \item Bereken die snypunte van $f$ en $g$.
    \item Gebruik dan jou grafieke en hulle snypunte om vir $x$ op te los wanneer:
	\begin{enumerate}[noitemsep, label=\textbf{\roman*}. ] 
	\item $f(x)>0$
	\item $\frac{f(x)}{g(x)}\le 0$
	\end{enumerate}
    \item Gee die vergelyking van die refleksie van $f$ in die $x$-as.
    \end{enumerate}
\item Nadat ’n bal neergegooi is, is die hoogte wat die bal terugbons elke keer minder. Die vergelyking $y=5{(0,8)}^{x}$ toon die verwantskap tussen
 $x$,die nommer van die bons, en $y$, die hoogte van die bons vir ’n
spesifieke bal. Wat is die benaderde hoogte van die vyfde bons tot die naaste tiende van ’n eenheid?\newline
\item Mark he $15$ muntstukke in R$~5$ en R$~2$ stukke. Hy het $3$ meer R$~2$ stukke as R$~5$ stukke. Hy het ‘n stelsel
van vergelykings opgestel om die situasie te toon, waar $x$  die hoeveelheid R$~5$ stukke voorstel en $y$ die
hoeveelheid R$~2$ stukke voorstel. Hy het vervolgens die probleem grafies opgelos.
    \begin{enumerate}[noitemsep, label=\textbf{\alph*}. ] 
    \item Skryf die sisteem van vergelykings neer.
    \item Skets die grafieke op dieselfde assestelsel.
    \item Wat is die oplossing?
    \end{enumerate}

\end{enumerate}

\par \raisebox{-0.2em}{\includegraphics[height=1em]{../icons/www.pdf}} Find the answers with the shortcodes:
\par \begin{tabular}[h]{cccccc}
(1.) lTN  &  (2.) lTR  &  (3.) lTp  &  (4.) lTn  &  (5.) lTy  &  (6.) lx8  &  (7.) lxX  &  (8.) lx9  & \end{tabular}
\end{eocexercises}